% Options for packages loaded elsewhere
\PassOptionsToPackage{unicode}{hyperref}
\PassOptionsToPackage{hyphens}{url}
\documentclass[
]{article}
\usepackage{xcolor}
\usepackage{amsmath,amssymb}
\setcounter{secnumdepth}{5}
\usepackage{iftex}
\ifPDFTeX
  \usepackage[T1]{fontenc}
  \usepackage[utf8]{inputenc}
  \usepackage{textcomp} % provide euro and other symbols
\else % if luatex or xetex
  \usepackage{unicode-math} % this also loads fontspec
  \defaultfontfeatures{Scale=MatchLowercase}
  \defaultfontfeatures[\rmfamily]{Ligatures=TeX,Scale=1}
\fi
\usepackage{lmodern}
\ifPDFTeX\else
  % xetex/luatex font selection
\fi
% Use upquote if available, for straight quotes in verbatim environments
\IfFileExists{upquote.sty}{\usepackage{upquote}}{}
\IfFileExists{microtype.sty}{% use microtype if available
  \usepackage[]{microtype}
  \UseMicrotypeSet[protrusion]{basicmath} % disable protrusion for tt fonts
}{}
\makeatletter
\@ifundefined{KOMAClassName}{% if non-KOMA class
  \IfFileExists{parskip.sty}{%
    \usepackage{parskip}
  }{% else
    \setlength{\parindent}{0pt}
    \setlength{\parskip}{6pt plus 2pt minus 1pt}}
}{% if KOMA class
  \KOMAoptions{parskip=half}}
\makeatother
\usepackage{longtable,booktabs,array}
\usepackage{caption}
\captionsetup[table]{skip=6pt}
\newcounter{none} % for unnumbered tables
\usepackage{calc} % for calculating minipage widths
% Correct order of tables after \paragraph or \subparagraph
\usepackage{etoolbox}
\makeatletter
\patchcmd\longtable{\par}{\if@noskipsec\mbox{}\fi\par}{}{}
\makeatother
% Allow footnotes in longtable head/foot
\IfFileExists{footnotehyper.sty}{\usepackage{footnotehyper}}{\usepackage{footnote}}
\makesavenoteenv{longtable}
\usepackage{graphicx}
\makeatletter
\newsavebox\pandoc@box
\newcommand*\pandocbounded[1]{% scales image to fit in text height/width
  \sbox\pandoc@box{#1}%
  \Gscale@div\@tempa{\textheight}{\dimexpr\ht\pandoc@box+\dp\pandoc@box\relax}%
  \Gscale@div\@tempb{\linewidth}{\wd\pandoc@box}%
  \ifdim\@tempb\p@<\@tempa\p@\let\@tempa\@tempb\fi% select the smaller of both
  \ifdim\@tempa\p@<\p@\scalebox{\@tempa}{\usebox\pandoc@box}%
  \else\usebox{\pandoc@box}%
  \fi%
}
% Set default figure placement to htbp
\def\fps@figure{htbp}
\makeatother
% definitions for citeproc citations
\NewDocumentCommand\citeproctext{}{}
\NewDocumentCommand\citeproc{mm}{%
  \begingroup\def\citeproctext{#2}\cite{#1}\endgroup}
\makeatletter
 % allow citations to break across lines
 \let\@cite@ofmt\@firstofone
 % avoid brackets around text for \cite:
 \def\@biblabel#1{}
 \def\@cite#1#2{{#1\if@tempswa , #2\fi}}
\makeatother
\newlength{\cslhangindent}
\setlength{\cslhangindent}{1.5em}
\newlength{\csllabelwidth}
\setlength{\csllabelwidth}{3em}
\newenvironment{CSLReferences}[2] % #1 hanging-indent, #2 entry-spacing
 {\begin{list}{}{%
  \setlength{\itemindent}{0pt}
  \setlength{\leftmargin}{0pt}
  \setlength{\parsep}{0pt}
  % turn on hanging indent if param 1 is 1
  \ifodd #1
   \setlength{\leftmargin}{\cslhangindent}
   \setlength{\itemindent}{-1\cslhangindent}
  \fi
  % set entry spacing
  \setlength{\itemsep}{#2\baselineskip}}}
 {\end{list}}
\usepackage{calc}
\newcommand{\CSLBlock}[1]{\hfill\break\parbox[t]{\linewidth}{\strut\ignorespaces#1\strut}}
\newcommand{\CSLLeftMargin}[1]{\parbox[t]{\csllabelwidth}{\strut#1\strut}}
\newcommand{\CSLRightInline}[1]{\parbox[t]{\linewidth - \csllabelwidth}{\strut#1\strut}}
\newcommand{\CSLIndent}[1]{\hspace{\cslhangindent}#1}
\setlength{\emergencystretch}{3em} % prevent overfull lines
\providecommand{\tightlist}{%
  \setlength{\itemsep}{0pt}\setlength{\parskip}{0pt}}
% symbols used in the Markdown source that the default text font lacks
\usepackage{newunicodechar}
\newunicodechar{≥}{\ensuremath{\geq}}
\newunicodechar{≤}{\ensuremath{\leq}}
\newunicodechar{∞}{\ensuremath{\infty}}
\newunicodechar{−}{\ensuremath{-}}
\usepackage{bookmark}
\IfFileExists{xurl.sty}{\usepackage{xurl}}{} % add URL line breaks if available
\urlstyle{same}
% fallback for those not using the hyperref driver hyperxmp:
\makeatletter
\@ifundefined{xmpquote}{\newcommand{\xmpquote}[1]{#1}}{}
\makeatother
\hypersetup{
  pdftitle={COSMOS-AI: a physics-aware{,} uncertainty-calibrated digital twin for orbital debris analysis{,} with atmospheric re-entry prediction and satellite health monitoring},
  pdfauthor={{[}Author name(s){]}a; a{[}Affiliation, postal address, country{]}; Corresponding author: {[}name, e-mail{]}},
  hidelinks,
  pdfcreator={LaTeX via pandoc}}

\title{COSMOS-AI: a physics-aware, uncertainty-calibrated digital twin
for orbital debris analysis, with atmospheric re-entry prediction and
satellite health monitoring}
\author{{[}Author
name(s){]}\textsuperscript{a} \and \textsuperscript{a}{[}Affiliation,
postal address, country{]} \and Corresponding author: {[}name,
e-mail{]}}
\date{}

\begin{document}
\maketitle
\begin{abstract}
Orbital debris analysis must forecast how long objects stay in orbit and
how likely they are to collide, yet practice rests on uncalibrated
physics propagators. Prior lifetime studies used a few hundred to 1,800
objects, single-epoch predictions and only decayed objects; tabular
foundation models have been assessed on curated benchmarks, not censored
catalogues under drift. This study proposes an artificial-intelligence
architecture for the debris environment, each component measured against
the practice it replaces. COSMOS-AI is a digital twin holding one state
per object. Debris forecasts collision risk from conjunction messages;
re-entry predicts residual lifetime by a reliability-gated mixture of
experts, a learned gate weighting a Monte-Carlo decay prior against
2025--2026 tabular models (TabM; in-context foundation models TabICLv2
and Kumo-Tabular); health scores telemetry with Chronos-2 zero-shot;
quantiles receive Mondrian conformal and isotonic calibration. Four
couplings make this one twin, not three independent models:
breakup-family statistics feed the lifetime learner, per-object survival
sums into cloud forecasts, operational status selects the decay regime,
and health flags escalate high-risk conjunctions. Public data
throughout: 64,516 catalogued objects in two forward-in-time periods,
the European Space Agency collision-avoidance benchmark and NASA
SMAP/MSL telemetry. In 2016--2021 the foundation models cut re-entry
pinball 13 \% below the best gradient-boosted baseline and 44 \% below
the physics expert, though tied with it in 2011--2015. Conformal
intervals covered 72.6 \%--90.6 \% of decayed objects against 12.2
\%--17.1 \% for the operational ±20 \% rule, designed for the final days
before re-entry, and history-conditioned state estimation cut 72-day
over-prediction from 9.5-fold to 1.4-fold. Collision risk improved 37 \%
over persistence, yet none beat the latest-risk baseline on the official
metric; zero-shot telemetry detection did not beat simple detectors.
Three couplings were testable --- family evidence helped in one period,
cloud survival reached a tenth of a cloud, the status gain was not
significant --- and the fourth lacks public test data. Calibrated
learned fusion advances debris forecasting where the full distribution
matters, and covariate availability is a leakage channel. A dated,
hashed forecast enables blind verification.

\textbf{Keywords:} Tabular foundation models; In-context learning;
Mixture of experts; Conformal prediction; Digital twin; Space debris;
Atmospheric re-entry prediction; Space situational awareness
\end{abstract}

\section{Introduction}\label{sec:intro}

More than 30,000 tracked objects share low Earth orbit with the
satellites that operators must keep healthy, steer clear of debris and
eventually bring down. These three tasks are usually handled by separate
tools and teams. Satellite health monitoring looks for anomalies in
telemetry; conjunction assessment estimates the probability that two
objects collide; and re-entry prediction estimates when an object will
decay and what that implies for people on the ground. Yet the tasks
share state. A satellite that loses its ability to manoeuvre stops
station-keeping and begins to decay naturally; fragments of one breakup
share an origin, and so a range of area-to-mass ratios, that shapes how
long they stay in orbit; and how long fragments stay up determines the
debris environment that conjunction assessment faces.

A digital twin of the tracked population, with one state per object that
all three functions read and update, is the natural way to make these
links explicit. Digital twins have been proposed for space systems
mostly as architectures and reviews (\citeproc{ref-liu2024digital}{Liu
et al., 2024}), or for single functions such as orbit prediction
(\citeproc{ref-xu2025digital}{Xu et al., 2025}); we found no system that
couples health, debris and re-entry functions and measures what the
coupling contributes. Building one on public data meets a hard
constraint: there is no public dataset that follows the same object
through telemetry, conjunctions and re-entry. Telemetry benchmarks are
anonymised (\citeproc{ref-hundman2018}{Hundman et al., 2018};
\citeproc{ref-kotowski2026esaadb}{Kotowski et al., 2026};
\citeproc{ref-ruszczak2025opssat}{Ruszczak et al., 2025}), conjunction
messages are anonymised (\citeproc{ref-uriot2022spacecraft}{Uriot et
al., 2022}), and catalogue data carry no telemetry. A credible twin
therefore cannot be a joint model trained on unpaired data. It has to be
a set of modules, each validated on its own real data, joined by
couplings that are defined explicitly and tested wherever a real test
exists.

Re-entry prediction is the module where the gap to practice is largest.
Operational re-entry predictions are physics-based and concentrate on
the final weeks, with windows of about ±20 \% of the residual lifetime
that cover the true epoch roughly 90 \% of the time
(\citeproc{ref-pardini2018assessing}{Pardini and Anselmo, 2018}).
Machine-learning studies of re-entry and lifetime have used Space-Track
orbital histories of a few hundred to about 1,800 objects, point
predictions and, as far as we could verify, only decayed objects
(\citeproc{ref-jung2021recurrent}{Jung et al., 2021};
\citeproc{ref-lotti2025exploring}{Lotti et al., 2025};
\citeproc{ref-salmaso2023machine}{Salmaso et al., 2023}). A
catalogue-wide model that handles objects still in orbit, is evaluated
forward in time across solar-cycle phases and produces intervals with a
coverage guarantee was not found in the literature we screened (Section
\ref{sec:related}).

This paper presents COSMOS-AI, a digital-twin framework for satellite
health monitoring, orbital debris analysis and atmospheric re-entry risk
prediction, and evaluates it on real data only. Its objectives are:

\begin{enumerate}
\def\labelenumi{\arabic{enumi}.}
\tightlist
\item
  to build and validate a catalogue-wide re-entry module that fuses a
  physics expert with learned experts and calibrates its uncertainty
  without distributional assumptions;
\item
  to validate a debris module (conjunction-risk forecasting and
  breakup-family evidence) and a health module (telemetry anomaly
  detection) on public benchmarks;
\item
  to define the couplings between the modules and measure, on real data,
  what each coupling contributes.
\end{enumerate}

The contributions are:

\begin{itemize}
\tightlist
\item
  \textbf{A twin framework with tested couplings.} One state per object,
  three modules sharing one uncertainty principle (conformal calibration
  of alarms, intervals and probabilities), and four couplings, three of
  which are tested on real data: breakup-family evidence for re-entry
  prediction, breakup-cloud survival forecasting, and operational status
  as health evidence that selects the re-entry regime. The fourth
  (health × debris → collision response) has no public test data and is
  labelled a demonstration. The twin's forecasts for objects already in
  orbit condition each object's area-to-mass ratio on its catalogue
  history. They are scored against later catalogue decays in an
  as-of-date test (T1) on a window that had been analysed earlier in the
  project; a dated, hashed forecast is released for a blind test.
\item
  \textbf{A licence-clean, leakage-controlled re-entry benchmark} built
  from GCAT (\citeproc{ref-mcdowell2020gcat}{McDowell, 2020}) and
  CelesTrak (\citeproc{ref-celestrak2026satcat}{Kelso, 2026}): 64,516
  objects with low-Earth-orbit canonical orbits, administrative
  censoring at the forecast origin, and two independent five-year test
  periods.
\item
  \textbf{A reliability-gated re-entry model} that averages the quantile
  functions of a Monte-Carlo physics expert and a learned expert, with
  Mondrian conformal intervals. The physics expert uses orbit-averaged
  King-Hele drag, exact for eccentric orbits, with NRLMSIS 2.1
  densities, and area-to-mass ratios from design values, class priors
  or, for fragments, the NASA Standard Breakup Model. The gate is driven
  by the physics expert's own reliability descriptors.
\item
  \textbf{A documented outcome leak in catalogue covariates.} Whether a
  catalogued object has a published radar cross-section depends on how
  long it stayed in orbit. Using it broke interval coverage in the test
  period whose missingness pattern differed from training; we exclude it
  and provide a check for such leaks.
\item
  \textbf{An honest account of where learning helps and where it does
  not}, including results in which the proposed components do not beat
  simpler alternatives.
\end{itemize}

\section{Related work}\label{sec:related}

\textbf{Re-entry and lifetime prediction.} Semi-analytic decay theory
(\citeproc{ref-frey2019extension}{Frey et al., 2019};
\citeproc{ref-kinghele1987satellite}{King-Hele, 1987}) and numerical
propagators with empirical atmospheres underpin operational lifetime
tools such as ESA's DRAMA/OSCAR, whose predictions were compared with
981 observed lifetimes by Operti et al.
(\citeproc{ref-operti2017assessing}{2017}). Assessments of the IADC
re-entry campaigns show mean errors of about 10 \% of the remaining
lifetime and a ±20 \% window covering about 90 \% of epochs near
re-entry (\citeproc{ref-pardini2018assessing}{Pardini and Anselmo,
2018}). Learning-based studies predict the re-entry epoch from orbital
histories: Salmaso et al. (\citeproc{ref-salmaso2023machine}{2023})
trained a sequence model on 417 Space-Track objects with
tracking-and-impact messages and tested it on 15 IADC objects; Jung et
al. (\citeproc{ref-jung2021recurrent}{2021}) used a recurrent network on
TLE-derived altitude histories; physics-informed networks for lifetime
are emerging (\citeproc{ref-lotti2025exploring}{Lotti et al., 2025}). A
recent physics-based backtest reports calibrated lifetime distributions
for 1,283 objects but calibrates on the same backtest set
(\citeproc{ref-shambaugh2026validation}{Shambaugh, 2026}). Uncertainty
in thermospheric density, the main driver, has been modelled with
calibrated neural networks (\citeproc{ref-licata2022uncertainty}{Licata
and Mehta, 2022}).

\textbf{Debris analysis and conjunction risk.} The ESA Kelvins Collision
Avoidance Challenge released Conjunction Data Messages with the final
collision risk of each event (\citeproc{ref-uriot2022spacecraft}{Uriot
et al., 2022}, \citeproc{ref-uriot2021cacdataset}{2021}). The
latest-risk baseline scored 0.694 on the official metric and only 12 of
97 teams beat it; the best entry scored 0.556
(\citeproc{ref-uriot2022spacecraft}{Uriot et al., 2022}). Later entries
report scores on the same split (0.628 for Tulczyjew et al.
(\citeproc{ref-tulczyjew2021predicting}{2021})), and a pre-registered
re-evaluation found that no model beats the latest-risk forecast with
confidence (\citeproc{ref-vazquez2026latest}{Vazquez, 2026}). Debris
families have been studied for origin attribution
(\citeproc{ref-celletti2021reconnecting}{Celletti et al., 2021};
\citeproc{ref-wu2023investigation}{Wu and Rosengren, 2023}); the Kelvins
``Space debris: the origin'' data are simulated
(\citeproc{ref-efthymiopoulos2021origin}{Efthymiopoulos et al., 2021}).
We found no study that uses breakup-family information as a predictor of
fragment lifetime.

\textbf{Telemetry anomaly detection.} ESA-ADB
(\citeproc{ref-kotowski2026esaadb}{Kotowski et al., 2026}), OPSSAT-AD
(\citeproc{ref-ruszczak2025opssat}{Ruszczak et al., 2025}) and NASA
SMAP/MSL (\citeproc{ref-hundman2018}{Hundman et al., 2018}) are the
public spacecraft benchmarks; flawed metrics such as point adjustment
can make random detectors look strong
(\citeproc{ref-kim2022rigorous}{Kim et al., 2022};
\citeproc{ref-wu2021flawed}{Wu and Keogh, 2021}), so event-wise and
point-wise metrics without adjustment are used here.

\textbf{Uncertainty.} Conformal prediction gives finite-sample,
distribution-free coverage under exchangeability
(\citeproc{ref-angelopoulos2023gentle}{Angelopoulos and Bates, 2023};
\citeproc{ref-vovk2005algorithmic}{Vovk et al., 2005}); conformalised
quantile regression adapts interval width to the input
(\citeproc{ref-romano2019cqr}{Romano et al., 2019}); Mondrian conformal
prediction conditions coverage on groups
(\citeproc{ref-vovk2005algorithmic}{Vovk et al., 2005}). Averaging
quantile functions is a well-studied way to combine forecasts
(\citeproc{ref-lichtendahl2013average}{Lichtendahl et al., 2013}). In
our reading, conformal intervals have not been applied to re-entry or
lifetime prediction.

\section{Data}\label{sec:data}

All data are public and real; their provenance, licences, checksums and
limitations are listed in the supplementary dataset audit.

\textbf{Catalogue (re-entry and debris modules).} GCAT
(\citeproc{ref-mcdowell2020gcat}{McDowell, 2020}) (release 1.8.8,
CC-BY-4.0, retrieved 2026-10-27) lists for each catalogued object one
\emph{canonical orbit}, ``representative of the spacecraft's early
operations, after initial positioning'', with its epoch; design mass and
dimensions; the parent object (for fragments, the object that broke up);
and the end of its current phase with a status code that separates
natural decay from active de-orbit, landing or docking. The CelesTrak
SATCAT (\citeproc{ref-celestrak2026satcat}{Kelso, 2026}) adds
operational status. CelesTrak's space-weather file provides F10.7 and Ap
since 1957; only observed rows of days that ended before the prediction
time are used. The SATCAT's radar cross-section (RCS) is \emph{not} used
as a predictor, because its availability encodes the outcome. Values
were mostly never published for objects that decayed quickly, and none
are published for objects catalogued after about 2018. In the
origin-2016 training cohort, 75.9 \% of objects without an RCS decayed
within 30 days, against 17.2 \% of those with one. Adjusted for orbit,
class and decade, the odds ratio is 3.59 (95 \% CI 2.09--6.17). The
effect is quantified in Section \ref{sec:results}. The cohort holds
64,516 objects whose canonical orbit has perigee below 2000 km, of which
32,970 decayed naturally; the exclusions are counted in the build log.
Fig. \ref{fig:data-reentry} shows the cohort's orbits, its distribution
over time with the cohorts of the two forecast origins, the observed
lifetimes and the space-weather record.

\begin{figure}
\centering
\includegraphics[width=1\linewidth,height=\textheight,keepaspectratio,alt={The re-entry dataset. (a) GCAT canonical orbits of the cohort by object class; payloads are outside the primary population. The classes overlap heavily at these densities, so the panel shows where each class lies rather than its density; the largest class is drawn first. (b) Objects per year of their canonical-orbit epoch, with the training, tuning and test cohorts of the two forecast origins. (c) Time from the canonical-orbit epoch to natural decay, or to censoring for objects still in orbit at the snapshot or whose phase ended without natural decay (a controlled re-entry, a landing, a docking or another non-decay status code); the learning target is capped at five years. (d) F10.7 solar flux from the CelesTrak space-weather file; its 81-day mean is the model input, and the bars mark the two test periods.}]{artwork/Figure_1.png}
\caption{\textbf{The re-entry dataset.} (a) GCAT canonical orbits of the
cohort by object class; payloads are outside the primary population. The
classes overlap heavily at these densities, so the panel shows where
each class lies rather than its density; the largest class is drawn
first. (b) Objects per year of their canonical-orbit epoch, with the
training, tuning and test cohorts of the two forecast origins. (c) Time
from the canonical-orbit epoch to natural decay, or to censoring for
objects still in orbit at the snapshot or whose phase ended without
natural decay (a controlled re-entry, a landing, a docking or another
non-decay status code); the learning target is capped at five years. (d)
F10.7 solar flux from the CelesTrak space-weather file; its 81-day mean
is the model input, and the bars mark the two test
periods.}\label{fig:data-reentry}
\end{figure}

\textbf{Conjunctions (debris module).} The Kelvins Collision Avoidance
Challenge data (\citeproc{ref-uriot2021cacdataset}{Uriot et al., 2021})
contain anonymised CDMs received by ESA in 2015--2019. Following the
challenge, inputs are the CDMs released at least two days before closest
approach and the target is the risk (log10 collision probability) of the
final CDM. Training events are filtered to mirror the test inputs,
leaving 8,293 training events (66 high-risk, i.e.~final risk ≥ −6) and
the official 2,167 test events (150 high-risk). Fig.
\ref{fig:data-debris} shows the breakup clouds of the catalogue and the
conjunction messages.

\begin{figure}
\centering
\includegraphics[width=1\linewidth,height=\textheight,keepaspectratio,alt={Debris data. (a) Catalogued fragments in orbit over time for the four largest breakup clouds of the cohort (fragments with a GCAT canonical orbit with perigee below 2000 km). (b) Gabbard diagram of the largest cloud, from the canonical orbits of its fragments. (c) Risk sequences of the conjunction data messages of 75 training events of the ESA Kelvins challenge: 60 drawn at random from the events whose final risk is below −6 and 15 from the 66 high-risk events (fixed seed), so high-risk events are deliberately oversampled. The model inputs are the messages received at least two days before closest approach. (d) Final risk of the training events: most lie at the floor of −30, and few are high-risk.}]{artwork/Figure_2.png}
\caption{\textbf{Debris data.} (a) Catalogued fragments in orbit over
time for the four largest breakup clouds of the cohort (fragments with a
GCAT canonical orbit with perigee below 2000 km). (b) Gabbard diagram of
the largest cloud, from the canonical orbits of its fragments. (c) Risk
sequences of the conjunction data messages of 75 training events of the
ESA Kelvins challenge: 60 drawn at random from the events whose final
risk is below −6 and 15 from the 66 high-risk events (fixed seed), so
high-risk events are deliberately oversampled. The model inputs are the
messages received at least two days before closest approach. (d) Final
risk of the training events: most lie at the floor of −30, and few are
high-risk.}\label{fig:data-debris}
\end{figure}

\textbf{Telemetry (health module).} NASA SMAP/MSL
(\citeproc{ref-hundman2018}{Hundman et al., 2018}): 81 channels with 68
SMAP and 36 MSL labelled anomaly sequences (the label file lists channel
P-2 twice; the two rows are merged). The alarm layer of the health
module was developed and tested on ESA-ADB Missions 1 and 2
(\citeproc{ref-kotowski2026esaadb}{Kotowski et al., 2026}) in a
companion pre-registered study by the same authors, whose results are
summarised in Section \ref{sec:res-health}. Fig.
\ref{fig:data-telemetry} shows examples of both.

\begin{figure}
\centering
\includegraphics[width=1\linewidth,height=\textheight,keepaspectratio,alt={Telemetry data. (a--d) Two channels of each NASA spacecraft drawn at random (fixed seed): training series (grey), test series (black) and labelled anomaly sequences (shaded). (e) ESA-ADB Mission 1 around the first labelled anomaly of the test half, for the target channels it affects (values min--max scaled per channel within the window; shaded: labels). ESA-ADB is used in the companion study.}]{artwork/Figure_3.png}
\caption{\textbf{Telemetry data.} (a--d) Two channels of each NASA
spacecraft drawn at random (fixed seed): training series (grey), test
series (black) and labelled anomaly sequences (shaded). (e) ESA-ADB
Mission 1 around the first labelled anomaly of the test half, for the
target channels it affects (values min--max scaled per channel within
the window; shaded: labels). ESA-ADB is used in the companion
study.}\label{fig:data-telemetry}
\end{figure}

\textbf{What is not available.} Orbital-element histories at catalogue
scale require a Space-Track account whose user agreement forbids
redistribution, and no legitimate public archive exists; DISCOS and ESA
re-entry predictions need accounts. The re-entry module therefore
predicts from a single dated orbit per object, not from orbital
histories.

\section{The COSMOS-AI framework}\label{sec:framework}

\subsection{Twin state}\label{twin-state}

For object \emph{i} at time \emph{t} the twin holds
\(s_i(t) = (o_i, \pi_i, h_i, F_i, e, C_i)\): the latest orbit \(o_i\)
(perigee, apogee, inclination), design properties \(\pi_i\), health
evidence \(h_i\) (catalogue operational status for payloads; the health
module's alarm state where telemetry exists), the debris family \(F_i\)
with sibling statistics restricted to what is known at \emph{t}, space
weather \(e\), and a conjunction summary \(C_i\) when CDMs exist. Every
module output attached to the twin records its model, validation status
and uncertainty (Fig. \ref{fig:arch}).

\begin{figure}
\centering
\includegraphics[width=1\linewidth,height=\textheight,keepaspectratio,alt={COSMOS-AI digital twin. Three modules read and update one state per object. Solid arrows are couplings tested on real data in this paper; the dashed coupling has no public test data.}]{artwork/Figure_4.png}
\caption{\textbf{COSMOS-AI digital twin.} Three modules read and update
one state per object. Solid arrows are couplings tested on real data in
this paper; the dashed coupling has no public test
data.}\label{fig:arch}
\end{figure}

\subsection{Re-entry module}\label{re-entry-module}

\emph{Physics expert.} The semi-major axis \(a\) and eccentricity \(e\)
are integrated with King-Hele's orbit-averaged drag equations for an
atmosphere locally exponential about perigee
(\citeproc{ref-kinghele1987satellite}{King-Hele, 1987}):

\begin{equation}\protect\phantomsection\label{eq:kh}{\Delta a_{\text{rev}} = -2\pi\delta a^2\rho_p e^{-c}\left[I_0 + 2eI_1 + \tfrac{3}{4}e^2(I_0+I_2)\right], \qquad c = ae/H,}\end{equation}

with \(I_n(c)\) modified Bessel functions. \(\rho_p\) and \(H\) are the
density and scale height at perigee from NRLMSIS 2.1
(\citeproc{ref-emmert2021nrlmsis}{Emmert et al., 2021}), with the
trailing 81-day F10.7 and Ap at \(t_0\) held constant.
\(\delta = F\,C_D A/m\) is the ballistic coefficient times the
atmospheric-rotation factor, and re-entry is declared when the perigee
falls below 120 km.

The series (shown truncated) holds only for small eccentricity. Its
perigee rate is 0.68 of the exact value at \(e = 0.2\) and changes sign
near \(e = 0.3\). For \(e > 0.05\) we therefore integrate the same orbit
averages exactly over the eccentric anomaly \(E\):
\(\Delta a_{\text{rev}} = -\delta\rho_p a^2 \int_0^{2\pi} e^{c(\cos E - 1)} (1+e\cos E)^{3/2}(1-e\cos E)^{-1/2}\,dE\).
The trapezoid rule is spectrally accurate for this periodic integrand;
its relative error is below \(10^{-12}\) up to \(c = 3000\).

The area-to-mass ratio is sampled per object from one of three sources:

\begin{itemize}
\tightlist
\item
  the design value (log-normal, 0.25 dex);
\item
  for intact payloads and rocket stages without design values, a design
  value drawn from the training-set distribution of their class;
\item
  for fragments and components without design values, the NASA Standard
  Breakup Model's bimodal distribution
  (\citeproc{ref-johnson2001nasa}{Johnson et al., 2001}), with a
  characteristic length drawn from the training-set distribution of
  radar-derived sizes of their class.
\end{itemize}

Each object has its own random stream, so its predictive distribution
does not depend on the other objects in a batch. Thirty-two Monte-Carlo
samples give a predictive distribution with quantiles
\(q^{\text{phys}}_\tau\).

\emph{Learned experts.} Three recent models predict quantiles
\(q^{\text{ml}}_\tau\) of \(z = \log_{10}\min(T, 5\,\text{yr})\):

\begin{itemize}
\tightlist
\item
  TabM, a parameter-efficient ensemble of MLPs
  (\citeproc{ref-gorishniy2025tabm}{Gorishniy et al., 2025});
\item
  TabICLv2 (\citeproc{ref-qu2026tabiclv2}{Qu et al., 2026}), an
  in-context tabular foundation model;
\item
  Kumo-Tabular (\citeproc{ref-nvidia2026kumo}{NVIDIA, 2026}), also an
  in-context tabular foundation model.
\end{itemize}

Their inputs are the orbit, design properties, space weather, family
statistics and summaries of the physics distribution. TabM gets a
non-crossing quantile head and the same search budget as every tuned
learner. The two foundation models are used with their default inference
settings. Kumo-Tabular's context is capped at 16,000 training rows by
device memory. The candidate with the lowest tuning loss becomes the
expert.

Gradient-boosted trees (\citeproc{ref-chen2016xgboost}{Chen and
Guestrin, 2016}; \citeproc{ref-ke2017lightgbm}{Ke et al., 2017};
\citeproc{ref-prokhorenkova2018catboost}{Prokhorenkova et al., 2018})
and a plain non-crossing quantile MLP are fitted and reported as classic
reference baselines; they are not eligible as the expert. A
censoring-aware accelerated-failure-time model (XGBoost AFT), trained on
all objects including those still in orbit, is a further reference
baseline.

\emph{Reliability-gated fusion.} With reliability descriptors \(r_i\) of
the physics expert (whether its area-to-mass ratio comes from design
values or a class prior, the spread of its own distribution,
eccentricity, perigee band, object class),

\begin{equation}\protect\phantomsection\label{eq:gate}{q_\tau(i) = g(r_i)\,q^{\text{phys}}_\tau(i) + \bigl(1-g(r_i)\bigr)\,q^{\text{ml}}_\tau(i), \qquad g(r) = \sigma(w^\top r + b).}\end{equation}

A convex combination of two quantile functions is a quantile function
(\citeproc{ref-lichtendahl2013average}{Lichtendahl et al., 2013}).
\((w, b)\) minimise the pinball loss on out-of-fold predictions of the
training period.

\emph{Mondrian conformal calibration.} On a calibration split disjoint
from fitting and gating, the scores
\(S = \max(q_{\text{lo}} - z, z - q_{\text{hi}})\) of each group \(k\)
(object class × perigee band, with fallback to class and pooled) give
the correction \(\hat Q_k\), the \(\lceil(1-\alpha)(n_k+1)\rceil\)-th
smallest score, and the interval
\([q_{\text{lo}} - \hat Q_k,\; q_{\text{hi}} + \hat Q_k]\)
(\citeproc{ref-romano2019cqr}{Romano et al., 2019};
\citeproc{ref-vovk2005algorithmic}{Vovk et al., 2005}). Probabilities
\(P(T \le H)\) are read from the fused quantile function and calibrated
on the same split by isotonic regression
(\citeproc{ref-zadrozny2002transforming}{Zadrozny and Elkan, 2002}).
Rank metrics are computed on the raw probabilities, because the isotonic
map ties scores; Brier score and ECE are computed on the calibrated
ones. The target has a point mass at the 5-year cap, so predicted
quantiles within \(10^{-6}\) of the cap (in log10 days) are set to it
exactly.

\emph{Twin state estimation for objects already in orbit.} The re-entry
module is fitted and validated at canonical-orbit epochs. A twin,
however, forecasts objects that have been in orbit for years. Their
survival since the canonical orbit selects low area-to-mass ratios,
which neither the class priors nor the learned experts know about.

For each fragment, rocket stage and component, the twin therefore
conditions the area-to-mass ratio on the object's catalogue history, a
coarse two-orbit form of ballistic-coefficient estimation from orbital
histories (\citeproc{ref-gondelach2017ballistic}{Gondelach et al.,
2017}; \citeproc{ref-pardini2008impact}{Pardini and Anselmo, 2008};
\citeproc{ref-sang2013estimation}{Sang et al., 2013}):

\begin{enumerate}
\def\labelenumi{\arabic{enumi}.}
\tightlist
\item
  128 prior samples are propagated from the GCAT canonical orbit to the
  as-of date with the realised trailing space weather.
\item
  Samples that would already have re-entered get weight zero.
\item
  The others are weighted by a Gaussian likelihood of the current
  perigee and apogee, with \(\sigma = 10\,\text{km} + 0.3\,|\Delta|\),
  where \(\Delta\) is the observed change since the canonical orbit.
\item
  32 posterior samples are drawn by systematic resampling. If the
  effective sample size is below 4, survival alone is used; if no sample
  survives, the 8 longest survivors are used.
\end{enumerate}

The physics expert then propagates the current orbit with the posterior
samples. Payloads keep the class priors, because their past manoeuvres
make the decay history uninformative.

\subsection{Debris module}\label{debris-module}

\emph{Family evidence.} For the launch family and, for fragments, the
breakup family of object \emph{i}: the number of siblings catalogued by
\(t_0\), the number and fraction that decayed naturally before \(t_0\),
and their mean log-lifetime. Only events before \(t_0\) enter.

\emph{Conjunction risk.} A persistence expert (the latest risk plus the
training distribution of final-minus-latest changes in its risk bin) and
the same learned experts (TabM, TabICLv2, Kumo-Tabular, with the classic
baselines alongside) predict quantiles of the final risk; they are fused
by the gate of Eq. \ref{eq:gate} and calibrated as above. Because the
official metric scores a point forecast, a calibrated P(high risk) also
gates the persistence point (predict the latest risk if P(high) exceeds
a threshold chosen on out-of-fold training data, otherwise −6.001).

\subsection{Health module}\label{health-module}

Per-channel non-conformity scores from a forecaster fitted on nominal
data are turned into conformal p-values against the channel's nominal
pool, aggregated across channels with Benjamini--Hochberg control and
persistence into alarm runs, as in the companion pre-registered ESA-ADB
study. For SMAP/MSL, the level and persistence are fixed in advance and
the detectors are:

\begin{itemize}
\tightlist
\item
  the Chronos-2 time-series foundation model
  (\citeproc{ref-ansari2025chronos2}{Ansari et al., 2025}), used
  zero-shot: rolling 8-step forecasts of the telemetry value from up to
  512 past steps, scored by the absolute residual of the forecast
  median;
\item
  three classic baselines: ridge autoregression, PCA reconstruction, and
  nearest-neighbour window distance;
\item
  the scores published with telemanom
  (\citeproc{ref-hundman2018}{Hundman et al., 2018}).
\end{itemize}

Each channel's alarms come from conformal p-values against its own
training-series scores:

\begin{itemize}
\tightlist
\item
  Chronos-2 is never fitted to these series, so its training scores are
  out of sample.
\item
  For the fitted baselines, the calibration scores are computed out of
  fold, in contiguous blocks with a purge gap.
\item
  An alarm is raised causally, at the third consecutive step with
  \(p \le \max(10^{-3}, 1/(n_\text{cal}+1))\).
\end{itemize}

\subsection{Couplings}\label{couplings}

\begin{itemize}
\tightlist
\item
  \textbf{C1 debris → re-entry:} family statistics are inputs of the
  re-entry experts.
\item
  \textbf{C2 re-entry → debris environment:} the expected number of a
  breakup cloud's fragments in orbit at time \emph{t} is
  \(\sum_{i:\,t_{0,i}\le t} S_i(t - t_{0,i})\), with \(S_i\) the fused
  survival function. \(S_i\) is read in two ways:

  \begin{itemize}
  \tightlist
  \item
    raw, from the quantile function;
  \item
    calibrated, by isotonic maps of \(P(T \le a)\) fitted on tune-B
    fragments at 25 ages and interpolated in log-age.
  \end{itemize}
\item
  \textbf{C3 health → re-entry:} a payload reported operational is in an
  operator-controlled regime, where drag physics does not determine
  re-entry. The twin withholds the drag forecast and uses the payload
  base rate of the preceding window. Non-operational payloads, rocket
  stages, components and fragments are in the natural-decay regime, and
  payloads with no reported status form a separate unknown regime.
\item
  \textbf{C4 health × debris → response:} a high-risk conjunction for a
  payload with an active health alarm or non-operational status is
  escalated because its manoeuvre capability is uncertain (demonstration
  only).
\end{itemize}

\section{Experimental protocol}\label{sec:protocol}

\textbf{Re-entry.} Prediction time \(t_0\) is the canonical-orbit epoch.
The event is natural decay; controlled endings are censored at their
date and objects in orbit at the snapshot. Under the protocol \emph{as
of origin O}, the training cohort has \(t_0 < O - 7\) yr, the tuning
cohort \(O - 7 \le t_0 < O - 5\) yr (split by launch family into tune-A
for early stopping, selection and checks, and tune-B for calibration),
and all their outcomes are censored at \emph{O}; and the test cohort
starts at \(O\) and ends where the five-year follow-up is complete at
the snapshot, or at the next origin: \$t\_0 \textless{} \$ 2016-01-01 at
origin 2011 (2011--2015) and \$t\_0 \textless{} \$ 2021-10-27 at origin
2016 (2016--2021). The two origins, 2011-01-01 and 2016-01-01, therefore
give two disjoint test periods. The primary population is objects that
cannot manoeuvre (fragments, rocket stages, components); payloads are
excluded from it because many are operator-controlled. Each tuned
learner receives the same budget of 12 random hyperparameter
configurations scored on tune-A, and every learner five seeds; TabICLv2
and Kumo-Tabular are used with their defaults. Metrics:

\begin{itemize}
\tightlist
\item
  mean pinball loss over seven quantiles;
\item
  MAE of the median in log10 days;
\item
  absolute error in days for objects that decayed within five years;
\item
  coverage, width and Winkler score of 80 \% and 90 \% intervals;
\item
  Harrell's C on all test objects
  (\citeproc{ref-harrell1982evaluating}{Harrell et al., 1982});
\item
  AUC-PR, AUC-ROC, Brier score and ECE for re-entry within 30 days, six
  months, one year and five years.
\end{itemize}

Confidence intervals come from 2,000 cluster-bootstrap replicates over
launch families, paired for model differences. Learner predictions are
cached under a fingerprint of the data, features and settings, so a
stale result cannot enter a comparison. As a labelled leak demonstration
(never a candidate), the selected learner is also refitted with the
radar cross-section added.

\textbf{Debris.} The official Kelvins split is used.

\begin{itemize}
\tightlist
\item
  Training events are split into five folds. For each fold, the tuned
  learners select their configuration by random search on the other four
  folds; the selected model then predicts the fold.
\item
  These nested out-of-fold predictions drive learner selection, the
  gate, isotonic calibration, the decision threshold and the conformal
  calibration.
\item
  The test models are tuned on a held-out fold of the training events
  and averaged over five seeds.
\item
  TabICLv2 and Kumo-Tabular have no search: their out-of-fold models use
  the other four folds as context, and their test models use all
  training events. Their out-of-fold losses therefore come from more
  training data than those of the tuned learners, which hold out one of
  the four folds for their search; this favours them in selection.
\item
  Predicted quantiles within \(10^{-4}\) of the risk floor (−30) are set
  to it exactly, because most final risks lie at the floor and the
  float32 spacing there is about \(2\times10^{-6}\).
\end{itemize}

The official metric is \(L = \text{MSE}_{\text{HR}}/F_2\) with low-risk
predictions clipped at −6.001 (\citeproc{ref-uriot2022spacecraft}{Uriot
et al., 2022}). We also report:

\begin{itemize}
\tightlist
\item
  pinball loss;
\item
  AUC-PR and AUC-ROC of the raw P(high risk);
\item
  the Brier score and ECE of its isotonic-calibrated value;
\item
  interval coverage.
\end{itemize}

\textbf{Health.} ESA-ADB as in the companion pre-registration (official
split and corrected event-wise \(F_{0.5}\)); SMAP/MSL with thresholds
fixed in advance and no test labels used.

\textbf{Couplings.} C1: ablation of the family features. C2: predicted
against catalogue counts of fragments in orbit, for every breakup cloud
with at least 30 fragments in a test cohort, at 1--5 years after its
first catalogued fragment. T1 (twin re-entry forecasts on current
orbits): the twin is run as of 2026-08-16, the data date of a CelesTrak
snapshot, and its P(re-entry within 30 and 72 days) is scored against
the decays recorded in the 2026-10-27 snapshot. Three variants are
compared for fragments, rocket stages and components:

\begin{itemize}
\tightlist
\item
  the lifetime pipeline, refitted as of 2026-08-16;
\item
  the physics expert with class priors;
\item
  the history-conditioned physics expert, which is the twin's choice.
\end{itemize}

All three were specified before the test, but the window is not fresh
(see \emph{Deviations and prior use of test data} below). T1 is
therefore an as-of-date test on a previously analysed window, not a
blind prospective test. The hashed forecast of 2026-10-27 is the blind
test.

C3: as-of-date test on the same previously analysed window, with the
CelesTrak operational status of 2026-08-16 and decays up to 2026-10-27,
for Earth-orbiting payloads with perigee below 600 km.

\begin{itemize}
\tightlist
\item
  \emph{As-of-date} policies use only information available on
  2026-08-16:

  \begin{itemize}
  \tightlist
  \item
    physics for every payload;
  \item
    the previous 72-day payload base rate for every payload;
  \item
    the status-aware twin (base rate for operational payloads, physics
    otherwise).
  \end{itemize}
\item
  Platt recalibrations, status-blind and status-specific, are fitted on
  the follow-up window itself and cross-fitted by launch. They are
  reported only as an in-window diagnostic of attainable calibration.
\item
  Confidence intervals resample launches.
\end{itemize}

\textbf{Deviations and prior use of test data.} The test sets of this
study were not untouched when the reported results were produced, and we
list what had been seen.

\begin{itemize}
\tightlist
\item
  \emph{Learner pool.} The first complete re-entry run (run 1;
  Supplementary Table S1) chose the expert among XGBoost, LightGBM,
  CatBoost, the plain MLP and TabICLv2. After its test results were
  known at both origins, and after the collision-risk test set had been
  scored with the same pool, we restricted the eligible pool of both the
  re-entry and the collision-risk experts to learners published
  2022--2026 (TabM, TabICLv2, Kumo-Tabular), so that the expert would be
  chosen among current methods; the classic learners became reference
  baselines. The selection rule, budget, seeds, features, gate and
  calibration were not changed. The reported run is run 2. Run 1
  selected the MLP at both origins, and its 2011 gated model was much
  worse than run 2's (Section \ref{sec:results}), so the 2011 gated
  result depends on this change.
\item
  \emph{Smoke tests.} Before run 2, engineering tests of the new
  learners printed two test-set scores: Kumo-Tabular's origin-2016
  pinball loss (0.051) and TabM's origin-2011 pinball loss (0.0729). No
  setting was chosen with them.
\item
  \emph{Repeated scoring.} An earlier version of every benchmark was
  scored on its test set and then corrected after an internal review;
  the superseded results are archived with the code. The re-entry and
  collision-risk test sets were therefore scored more than once during
  development.
\item
  \emph{The T1 and C3 window.} An earlier version of the project fitted
  and tested re-entry classifiers (physics, logistic regression,
  XGBoost, XGBoost-AFT and a hazard network) on the window from
  2026-08-16 to 2026-10-23, and the aggregate decay count of the window
  was seen while the lifetime pipeline's bias on current orbits was
  diagnosed, before the history-conditioned variant was designed. A
  defect in the twin's payload prior (payloads without design values
  fell back to the fragment breakup law) was found and fixed after T1
  had first been scored; the fix changes payload probabilities only.
\item
  \emph{Post-hoc analyses.} The class breakdown, the gate applied to
  every learner, the head-to-head comparisons, the calibrated C2 reading
  and the C2 bootstrap were added after test results had been seen. They
  are labelled as post hoc, and all intervals are nominal 95 \%
  intervals without correction for multiplicity.
\end{itemize}

\section{Results}\label{sec:results}

\subsection{Re-entry prediction at canonical-orbit
epochs}\label{sec:res-reentry}

Table \ref{tbl:reentry} lists every re-entry model on both test periods
and Fig. \ref{fig:reentry} shows their pinball losses.

\begin{longtable}[]{@{}
  >{\raggedright\arraybackslash}p{(\linewidth - 20\tabcolsep) * \real{0.0909}}
  >{\raggedright\arraybackslash}p{(\linewidth - 20\tabcolsep) * \real{0.0909}}
  >{\raggedright\arraybackslash}p{(\linewidth - 20\tabcolsep) * \real{0.0909}}
  >{\raggedright\arraybackslash}p{(\linewidth - 20\tabcolsep) * \real{0.0909}}
  >{\raggedright\arraybackslash}p{(\linewidth - 20\tabcolsep) * \real{0.0909}}
  >{\raggedright\arraybackslash}p{(\linewidth - 20\tabcolsep) * \real{0.0909}}
  >{\raggedright\arraybackslash}p{(\linewidth - 20\tabcolsep) * \real{0.0909}}
  >{\raggedright\arraybackslash}p{(\linewidth - 20\tabcolsep) * \real{0.0909}}
  >{\raggedright\arraybackslash}p{(\linewidth - 20\tabcolsep) * \real{0.0909}}
  >{\raggedright\arraybackslash}p{(\linewidth - 20\tabcolsep) * \real{0.0909}}
  >{\raggedright\arraybackslash}p{(\linewidth - 20\tabcolsep) * \real{0.0909}}@{}}
\caption{\textbf{Re-entry prediction on the two test periods} (left:
origin 2016, test 2016--2021; right: origin 2011, test 2011--2015).
Pinball loss over seven quantiles and MAE of the median, both in log10
days; coverage and mean width (log10 days) of the Mondrian conformal 90
\% interval; Harrell's C. Models marked ``baseline'' are not eligible as
the gated expert.}\label{tbl:reentry}\tabularnewline
\toprule\noalign{}
\begin{minipage}[b]{\linewidth}\raggedright
Model
\end{minipage} & \begin{minipage}[b]{\linewidth}\raggedright
Pinball
\end{minipage} & \begin{minipage}[b]{\linewidth}\raggedright
MAE
\end{minipage} & \begin{minipage}[b]{\linewidth}\raggedright
Cov. 90 \%
\end{minipage} & \begin{minipage}[b]{\linewidth}\raggedright
Width 90 \%
\end{minipage} & \begin{minipage}[b]{\linewidth}\raggedright
C
\end{minipage} & \begin{minipage}[b]{\linewidth}\raggedright
Pinball
\end{minipage} & \begin{minipage}[b]{\linewidth}\raggedright
MAE
\end{minipage} & \begin{minipage}[b]{\linewidth}\raggedright
Cov. 90 \%
\end{minipage} & \begin{minipage}[b]{\linewidth}\raggedright
Width 90 \%
\end{minipage} & \begin{minipage}[b]{\linewidth}\raggedright
C
\end{minipage} \\
\midrule\noalign{}
\endfirsthead
\toprule\noalign{}
\begin{minipage}[b]{\linewidth}\raggedright
Model
\end{minipage} & \begin{minipage}[b]{\linewidth}\raggedright
Pinball
\end{minipage} & \begin{minipage}[b]{\linewidth}\raggedright
MAE
\end{minipage} & \begin{minipage}[b]{\linewidth}\raggedright
Cov. 90 \%
\end{minipage} & \begin{minipage}[b]{\linewidth}\raggedright
Width 90 \%
\end{minipage} & \begin{minipage}[b]{\linewidth}\raggedright
C
\end{minipage} & \begin{minipage}[b]{\linewidth}\raggedright
Pinball
\end{minipage} & \begin{minipage}[b]{\linewidth}\raggedright
MAE
\end{minipage} & \begin{minipage}[b]{\linewidth}\raggedright
Cov. 90 \%
\end{minipage} & \begin{minipage}[b]{\linewidth}\raggedright
Width 90 \%
\end{minipage} & \begin{minipage}[b]{\linewidth}\raggedright
C
\end{minipage} \\
\midrule\noalign{}
\endhead
\bottomrule\noalign{}
\endlastfoot
Physics expert alone & 0.0892 & 0.230 & 0.876 & 0.67 & 0.739 & 0.0527 &
0.148 & 0.974 & 1.03 & 0.766 \\
XGBoost (baseline) & 0.0576 & 0.190 & 0.921 & 0.80 & 0.888 & 0.0648 &
0.231 & 0.894 & 1.01 & 0.807 \\
LightGBM (baseline) & 0.0662 & 0.191 & 0.971 & 1.00 & 0.887 & 0.0683 &
0.221 & 0.981 & 1.18 & 0.811 \\
CatBoost (baseline) & 0.0611 & 0.207 & 0.911 & 0.93 & 0.878 & 0.0627 &
0.201 & 0.788 & 0.98 & 0.821 \\
Quantile MLP (baseline) & 0.0578 & 0.193 & 0.892 & 0.84 & 0.851 & 0.1147
& 0.390 & 0.852 & 1.40 & 0.747 \\
XGBoost AFT, censoring-aware (baseline) & 0.0657 & 0.185 & 0.901 & 0.62
& 0.866 & 0.0531 & 0.156 & 0.961 & 0.87 & 0.813 \\
TabM (2025) & 0.0593 & 0.192 & 0.877 & 0.68 & 0.861 & 0.0994 & 0.354 &
0.918 & 1.38 & 0.783 \\
TabICLv2 (2026) & 0.0499 & 0.166 & 0.962 & 1.00 & 0.894 & 0.0596 & 0.160
& 0.919 & 1.29 & 0.844 \\
Kumo-Tabular (2026) & 0.0508 & 0.162 & 0.974 & 1.07 & 0.899 & 0.0616 &
0.163 & 0.966 & 1.34 & 0.845 \\
Physics + expert, equal weights & 0.0691 & 0.205 & 0.838 & 0.64 & 0.859
& 0.0527 & 0.148 & 0.970 & 1.04 & 0.851 \\
Physics + expert, constant weight & 0.0601 & 0.193 & 0.877 & 0.67 &
0.861 & 0.0608 & 0.162 & 0.960 & 1.31 & 0.847 \\
Reliability-gated, gate fitted on tune-A & 0.0695 & 0.202 & 0.832 & 0.62
& 0.863 & 0.0616 & 0.163 & 0.959 & 1.34 & 0.845 \\
\textbf{Reliability-gated (COSMOS)} & 0.0604 & 0.193 & 0.874 & 0.67 &
0.860 & 0.0583 & 0.157 & 0.959 & 1.21 & 0.848 \\
\end{longtable}

\begin{figure}
\centering
\includegraphics[width=1\linewidth,height=\textheight,keepaspectratio,alt={Test pinball loss of every re-entry model on the two test periods. Grey: physics expert; blue: reference baselines, fixed-weight fusions and the variant whose gate is fitted on tune-A; green: the 2025--2026 learned experts; orange: the reliability-gated model.}]{artwork/Figure_5.png}
\caption{\textbf{Test pinball loss of every re-entry model} on the two
test periods. Grey: physics expert; blue: reference baselines,
fixed-weight fusions and the variant whose gate is fitted on tune-A;
green: the 2025--2026 learned experts; orange: the reliability-gated
model.}\label{fig:reentry}
\end{figure}

\subsubsection{Test period 2016--2021 (origin
2016)}\label{test-period-20162021-origin-2016}

\textbf{Selection.} On tune-A, TabM had the lowest loss among the recent
learners (Supplementary Fig. S1 plots tune-A against test loss for every
learner):

{\def\LTcaptype{none} % do not increment counter
\begin{longtable}[]{@{}ll@{}}
\toprule\noalign{}
Learner & Tune-A pinball \\
\midrule\noalign{}
\endhead
\bottomrule\noalign{}
\endlastfoot
TabM & 0.0283 \\
TabICLv2 & 0.0369 \\
Kumo-Tabular & 0.0433 \\
\end{longtable}
}

On the test cohort the ranking is different, and TabM falls behind both
foundation models (Table \ref{tbl:reentry}):

{\def\LTcaptype{none} % do not increment counter
\begin{longtable}[]{@{}ll@{}}
\toprule\noalign{}
Model & Test pinball \\
\midrule\noalign{}
\endhead
\bottomrule\noalign{}
\endlastfoot
TabICLv2 & 0.0499 \\
Kumo-Tabular & 0.0508 \\
TabM & 0.0593 \\
Classic baselines (best: XGBoost) & 0.0576 \\
Physics expert alone & 0.0892 \\
\end{longtable}
}

Both foundation models are significantly better than the selected TabM
(post-hoc paired comparisons, Supplementary Table S16; a negative
difference favours the first model):

{\def\LTcaptype{none} % do not increment counter
\begin{longtable}[]{@{}lll@{}}
\toprule\noalign{}
Comparison & Pinball difference & 95 \% CI \\
\midrule\noalign{}
\endhead
\bottomrule\noalign{}
\endlastfoot
TabICLv2 minus TabM & -0.0093 & -0.0162 to -0.0036 \\
Kumo-Tabular minus TabM & -0.0085 & -0.0167 to -0.0017 \\
TabICLv2 minus Kumo-Tabular & -0.0008 & -0.0036 to +0.0024 \\
\end{longtable}
}

The two foundation models cannot be separated from each other.

\textbf{Coverage and concordance.} The raw 90 \% quantile intervals of
the two in-context foundation models already cover 95.8 \% (TabICLv2)
and 97.3 \% (Kumo-Tabular) of test objects. TabM's cover 80.8 \%.

After Mondrian calibration on tune-B, TabM and the gated model remain
below nominal coverage (87.7 \% and 87.4 \%). This is consistent with a
loss of exchangeability between the 2009--2010 calibration cohort and
the test cohort, although the shortfall is small relative to the
launch-clustered sampling error. Their intervals are narrower, and their
Winkler scores slightly lower:

{\def\LTcaptype{none} % do not increment counter
\begin{longtable}[]{@{}ll@{}}
\toprule\noalign{}
Model & Winkler score (90 \%) \\
\midrule\noalign{}
\endhead
\bottomrule\noalign{}
\endlastfoot
TabM & 1.099 \\
TabICLv2 & 1.129 \\
\end{longtable}
}

Supplementary Fig. S2 plots coverage against interval width for every
model, and Fig. \ref{fig:reentry-outputs} shows individual forecasts
and, for every test object with a known five-year outcome, predicted
against observed outcomes.

The foundation models also have the best concordance (C = 0.894 and
0.899). A tuning cohort five to seven years older than the test cohort
did not identify the model that generalises best.

\textbf{Gated model.} The reliability-gated model fuses the physics
expert with TabM; its mean gate weight on physics on the test cohort is
0.11. Its pinball loss of 0.0604 compares as follows (paired cluster
bootstrap; a positive difference favours the gated model):

{\def\LTcaptype{none} % do not increment counter
\begin{longtable}[]{@{}lll@{}}
\toprule\noalign{}
Comparison & Pinball difference & 95 \% CI \\
\midrule\noalign{}
\endhead
\bottomrule\noalign{}
\endlastfoot
Physics expert alone & +0.0289 & +0.0148 to +0.0469 \\
TabM alone & -0.0011 & -0.0028 to +0.0002 \\
TabICLv2 & -0.0104 & -0.0179 to -0.0043 \\
Kumo-Tabular & -0.0096 & -0.0181 to -0.0023 \\
\end{longtable}
}

The gated model is clearly better than physics alone and no different
from TabM alone. It is significantly worse than both foundation models.

\textbf{Horizon probabilities.} The foundation models also lead for
re-entry within 30 days:

{\def\LTcaptype{none} % do not increment counter
\begin{longtable}[]{@{}lll@{}}
\toprule\noalign{}
Model & AUC-PR & Brier score \\
\midrule\noalign{}
\endhead
\bottomrule\noalign{}
\endlastfoot
Kumo-Tabular & 0.883 & 0.036 \\
TabICLv2 & 0.874 & 0.035 \\
Gated model & 0.829 & 0.051 \\
Physics expert alone & 0.719 & 0.051 \\
\end{longtable}
}

\textbf{Solar activity.} Solar cycle 25 rose faster than the persistence
assumption allows. Replaying the realised F10.7 lowers the physics
expert's pinball loss from 0.0892 to 0.0783. This is a diagnostic only:
solar activity after the prediction time accounts for a substantial part
of the physics expert's error in this period.

\begin{figure}
\centering
\includegraphics[width=1\linewidth,height=\textheight,keepaspectratio,alt={What the re-entry module outputs (origin 2016, test 2016--2021). (a) Forty test objects drawn at random (fixed seed) among those with a known five-year outcome, sorted by outcome: median and 90 \% Mondrian conformal interval of the physics expert and of the reliability-gated model, against the observed time to decay (diamonds) or survival beyond the five-year cap (triangles). (b--d) Predicted median against observed outcome for every such test object, for the physics expert, TabICLv2 and the gated model.}]{artwork/Figure_6.png}
\caption{\textbf{What the re-entry module outputs} (origin 2016, test
2016--2021). (a) Forty test objects drawn at random (fixed seed) among
those with a known five-year outcome, sorted by outcome: median and 90
\% Mondrian conformal interval of the physics expert and of the
reliability-gated model, against the observed time to decay (diamonds)
or survival beyond the five-year cap (triangles). (b--d) Predicted
median against observed outcome for every such test object, for the
physics expert, TabICLv2 and the gated
model.}\label{fig:reentry-outputs}
\end{figure}

\subsubsection{Test period 2011--2015 (origin
2011)}\label{test-period-20112015-origin-2011}

\textbf{Selection.} At this origin the tuning cohort is small: 248
objects in tune-A and 332 in tune-B. Kumo-Tabular had the lowest tune-A
loss:

{\def\LTcaptype{none} % do not increment counter
\begin{longtable}[]{@{}ll@{}}
\toprule\noalign{}
Learner & Tune-A pinball \\
\midrule\noalign{}
\endhead
\bottomrule\noalign{}
\endlastfoot
Kumo-Tabular & 0.0547 \\
TabM & 0.0561 \\
TabICLv2 & 0.0581 \\
\end{longtable}
}

\textbf{Physics expert against the foundation models.} The physics
expert alone has the lowest point estimate of the test pinball loss
(0.0527), but it cannot be separated from the foundation models or from
the gated model (post-hoc paired comparisons, Supplementary Table S16; a
negative difference favours the physics expert):

{\def\LTcaptype{none} % do not increment counter
\begin{longtable}[]{@{}lll@{}}
\toprule\noalign{}
Physics expert minus & Pinball difference & 95 \% CI \\
\midrule\noalign{}
\endhead
\bottomrule\noalign{}
\endlastfoot
TabICLv2 & -0.0069 & -0.0217 to +0.0122 \\
Kumo-Tabular & -0.0089 & -0.0253 to +0.0133 \\
Gated model & -0.0056 & -0.0199 to +0.0139 \\
\end{longtable}
}

The physics expert's 90 \% conformal coverage is 97.4 \%, but its
concordance is low (C = 0.766).

\textbf{TabM.} TabM, the expert selected at 2016, has a much higher loss
here (0.0994), as does the plain MLP (0.1147).

\textbf{By object class} (post hoc, Supplementary Table S2 and Fig. S3):

{\def\LTcaptype{none} % do not increment counter
\begin{longtable}[]{@{}llll@{}}
\toprule\noalign{}
Model & Fragments & Rocket stages & Components \\
\midrule\noalign{}
\endhead
\bottomrule\noalign{}
\endlastfoot
TabM & 0.1122 & 0.0473 & 0.0597 \\
Physics expert & 0.0432 & 0.0733 & 0.0959 \\
\end{longtable}
}

TabM's error is concentrated in fragments. The physics expert has the
lowest loss for fragments, while the learned experts have lower losses
for rocket stages and components (e.g.~TabICLv2: 0.0469 and 0.0703).

\textbf{Gated model.} The gated model fuses the physics expert with
Kumo-Tabular (mean physics weight 0.12) and scores 0.0583. Paired
differences (a positive difference favours the gated model):

{\def\LTcaptype{none} % do not increment counter
\begin{longtable}[]{@{}lll@{}}
\toprule\noalign{}
Comparison & Pinball difference & 95 \% CI \\
\midrule\noalign{}
\endhead
\bottomrule\noalign{}
\endlastfoot
Kumo-Tabular alone & +0.0033 & +0.0007 to +0.0055 \\
XGBoost & +0.0066 & +0.0017 to +0.0136 \\
LightGBM & +0.0101 & +0.0023 to +0.0215 \\
CatBoost & +0.0044 & -0.0011 to +0.0122 \\
TabM & +0.0412 & +0.0011 to +0.0905 \\
TabICLv2 & +0.0013 & -0.0016 to +0.0044 \\
Physics expert alone & -0.0056 & -0.0199 to +0.0139 \\
\end{longtable}
}

The 95 \% intervals exclude zero for Kumo-Tabular alone, XGBoost,
LightGBM and TabM, but not for CatBoost, TabICLv2 or the physics expert.
The censoring-aware AFT baseline scores 0.0531.

\textbf{Solar activity.} Replaying the realised F10.7 changes the
physics expert's pinball loss only from 0.0527 to 0.0516.

\subsubsection{Across both origins}\label{across-both-origins}

All intervals in this section are nominal 95 \% intervals without
correction for multiplicity. Apart from the comparisons of the gated
model, the analyses below were added after the test results were seen,
so we read them as exploratory. Fig. \ref{fig:reentry-benchmark}
collects the head-to-head comparisons.

\begin{enumerate}
\def\labelenumi{\arabic{enumi}.}
\tightlist
\item
  \textbf{TabICLv2 was best or tied in both periods.} At 2016 the two
  foundation models are significantly better than the physics expert,
  the gated model and the best tree baseline; for example, TabICLv2
  minus XGBoost is -0.0076 (-0.0121 to -0.0036). At 2011 TabICLv2, the
  physics expert and the gated model are statistically tied;
  Kumo-Tabular is tied with TabICLv2 and the physics expert but
  significantly worse than the gated model that fuses it (+0.0033,
  +0.0007 to +0.0055).
\item
  \textbf{Selection on the tuning cohort chose a significantly worse
  expert at 2016.} There, TabM was selected over the two foundation
  models. At 2011 the selected Kumo-Tabular is tied with TabICLv2.
\item
  \textbf{Gating improved the expert it fused at 2011 and was neutral at
  2016.} In a post-hoc sensitivity analysis that gates every learner
  (Supplementary Table S3):

  \begin{itemize}
  \tightlist
  \item
    at 2011 the intervals exclude zero for every learner to which the
    gate gave non-negligible weight (TabICLv2: +0.0027);
  \item
    at 2016 no interval excludes zero, and the point estimates are
    mostly slightly negative (TabICLv2 -0.0046, -0.0112 to +0.0007).
    This is when the physics expert's own error was largest.
  \end{itemize}
\end{enumerate}

\textbf{First learner set.} The deviation record (Section
\ref{sec:protocol}, Supplementary Table S1) used XGBoost, LightGBM,
CatBoost, the plain MLP and TabICLv2 as candidates. It selected the MLP
at both origins. Its gated model scored 0.1074 at 2011 and 0.0592 at
2016. The reported 2011 gated model therefore depends on the change of
the candidate pool.

\textbf{Operational windows} (Supplementary Table S5). The ±20 \%
residual-lifetime window around the physics median, the rule used for
the final days before re-entry
(\citeproc{ref-pardini2018assessing}{Pardini and Anselmo, 2018}), covers
12.2 \% (2016) and 17.1 \% (2011) of the test objects that decayed
within five years, against 72.6 \% and 90.6 \% for the physics expert's
90 \% conformal interval. The shortfall is a property of forecasts made
from a canonical orbit years earlier, not of the horizon: the window
covers a similar share of the objects that decayed within a week of
their canonical-orbit epoch. The rule's reported \textasciitilde90 \%
coverage refers to the final days before re-entry, when the orbit is
tracked continuously. Fig. \ref{fig:coverage} shows the same rules over
all test objects, including those still in orbit after five years, where
the ±20 \% window scores higher because it is wide enough to contain the
five-year cap.

\begin{figure}
\centering
\includegraphics[width=1\linewidth,height=\textheight,keepaspectratio,alt={Head-to-head re-entry comparisons (post hoc): differences in test pinball loss with paired launch-family cluster-bootstrap 95 \% intervals; a negative difference favours the first model. Filled markers: the nominal interval, uncorrected for multiplicity, excludes zero. ``Its expert'' is TabM at origin 2016 and Kumo-Tabular at origin 2011.}]{artwork/Figure_7.png}
\caption{\textbf{Head-to-head re-entry comparisons} (post hoc):
differences in test pinball loss with paired launch-family
cluster-bootstrap 95 \% intervals; a negative difference favours the
first model. Filled markers: the nominal interval, uncorrected for
multiplicity, excludes zero. ``Its expert'' is TabM at origin 2016 and
Kumo-Tabular at origin 2011.}\label{fig:reentry-benchmark}
\end{figure}

\begin{figure}
\centering
\includegraphics[width=1\linewidth,height=\textheight,keepaspectratio,alt={Interval coverage on the test periods, over all test objects, including those still in orbit after five years: the operational ±20 \% and ±30 \% residual-lifetime windows around the physics median, the raw quantile intervals of the gated model, and its marginal and Mondrian conformal intervals. Supplementary Table S5 gives the same rules restricted to the objects that decayed within five years.}]{artwork/Figure_8.png}
\caption{\textbf{Interval coverage on the test periods}, over all test
objects, including those still in orbit after five years: the
operational ±20 \% and ±30 \% residual-lifetime windows around the
physics median, the raw quantile intervals of the gated model, and its
marginal and Mondrian conformal intervals. Supplementary Table S5 gives
the same rules restricted to the objects that decayed within five
years.}\label{fig:coverage}
\end{figure}

\subsubsection{An outcome leak through catalogue
covariates}\label{sec:res-leak}

Whether an object has a radar cross-section in the 2026 catalogue
depends on its fate.

\textbf{Evidence of the leak.} In the origin-2016 training cohort:

\begin{itemize}
\tightlist
\item
  75.9 \% of objects without a value decayed within 30 days of their
  canonical-orbit epoch, against 17.2 \% of objects with one.
\item
  Adjusted for orbit, class and decade, the odds ratio is 3.59. Its 95
  \% CI, 2.09--6.17, uses cluster-robust standard errors over launch
  families.
\end{itemize}

The test periods differ in how this availability relates to the outcome:

\begin{itemize}
\tightlist
\item
  In the 2011--2015 test cohort, only 1.9 \% of objects lack a value,
  and 84.0 \% of them decayed within 30 days.
\item
  In the 2016--2021 cohort, the value is missing for 99.9 \% of objects
  catalogued from 2018 on, whatever their fate.
\end{itemize}

\textbf{Effect on the selected learner} (refitted with the 2026 RCS
added, a labelled demonstration and never a candidate):

{\def\LTcaptype{none} % do not increment counter
\begin{longtable}[]{@{}
  >{\raggedright\arraybackslash}p{(\linewidth - 10\tabcolsep) * \real{0.1667}}
  >{\raggedright\arraybackslash}p{(\linewidth - 10\tabcolsep) * \real{0.1667}}
  >{\raggedright\arraybackslash}p{(\linewidth - 10\tabcolsep) * \real{0.1667}}
  >{\raggedright\arraybackslash}p{(\linewidth - 10\tabcolsep) * \real{0.1667}}
  >{\raggedright\arraybackslash}p{(\linewidth - 10\tabcolsep) * \real{0.1667}}
  >{\raggedright\arraybackslash}p{(\linewidth - 10\tabcolsep) * \real{0.1667}}@{}}
\toprule\noalign{}
\begin{minipage}[b]{\linewidth}\raggedright
Origin
\end{minipage} & \begin{minipage}[b]{\linewidth}\raggedright
Pinball with RCS
\end{minipage} & \begin{minipage}[b]{\linewidth}\raggedright
Pinball without
\end{minipage} & \begin{minipage}[b]{\linewidth}\raggedright
Difference (95 \% CI)
\end{minipage} & \begin{minipage}[b]{\linewidth}\raggedright
90 \% coverage with RCS
\end{minipage} & \begin{minipage}[b]{\linewidth}\raggedright
90 \% coverage without
\end{minipage} \\
\midrule\noalign{}
\endhead
\bottomrule\noalign{}
\endlastfoot
2011 & 0.0578 & 0.0616 & -0.0038 (-0.0072 to +0.0011) & 96.1 \% & 96.6
\% \\
2016 & 0.0664 & 0.0593 & +0.0071 (-0.0010 to +0.0151) & 76.8 \% & 87.7
\% \\
\end{longtable}
}

\textbf{Pinball loss.} The change from adding the RCS is within noise at
both origins. It is also not consistent across learner sets: with the
first set, adding it worsened the selected MLP at 2011 (0.1221 against
0.1147).

\textbf{Interval coverage.} The consistent effect is on coverage in the
period whose missingness pattern differs from training. At 2016, 90 \%
coverage falls from 87.7 \% to 76.8 \%. With the first learner set it
falls from 89.2 \% to 72.1 \%.

\textbf{Conclusion.} The association between availability and outcome is
clear, and using the covariate breaks calibration when that association
changes. Its effect on point accuracy was not detectable here.

\subsubsection{Feature groups and coupling C1 (family
evidence)}\label{sec:res-c1}

Each feature group was removed from the selected learner and the learner
refitted with the same settings and seeds. Fig. \ref{fig:ablation} gives
the paired launch-family bootstrap of the change in test pinball loss; a
positive change means the group helps. The intervals are nominal 95 \%
intervals without correction for the eight comparisons.

{\def\LTcaptype{none} % do not increment counter
\begin{longtable}[]{@{}
  >{\raggedright\arraybackslash}p{(\linewidth - 8\tabcolsep) * \real{0.2000}}
  >{\raggedright\arraybackslash}p{(\linewidth - 8\tabcolsep) * \real{0.2000}}
  >{\raggedright\arraybackslash}p{(\linewidth - 8\tabcolsep) * \real{0.2000}}
  >{\raggedright\arraybackslash}p{(\linewidth - 8\tabcolsep) * \real{0.2000}}
  >{\raggedright\arraybackslash}p{(\linewidth - 8\tabcolsep) * \real{0.2000}}@{}}
\toprule\noalign{}
\begin{minipage}[b]{\linewidth}\raggedright
Group removed
\end{minipage} & \begin{minipage}[b]{\linewidth}\raggedright
Origin 2016 (TabM)
\end{minipage} & \begin{minipage}[b]{\linewidth}\raggedright
95 \% CI
\end{minipage} & \begin{minipage}[b]{\linewidth}\raggedright
Origin 2011 (Kumo-Tabular)
\end{minipage} & \begin{minipage}[b]{\linewidth}\raggedright
95 \% CI
\end{minipage} \\
\midrule\noalign{}
\endhead
\bottomrule\noalign{}
\endlastfoot
Family evidence (C1) & +0.0056 & +0.0015 to +0.0101 & +0.0024 & -0.0058
to +0.0103 \\
Physics summaries & -0.0050 & -0.0086 to -0.0022 & +0.0022 & +0.0002 to
+0.0039 \\
Design properties & +0.0071 & +0.0035 to +0.0121 & +0.0008 & -0.0001 to
+0.0017 \\
Space weather & +0.0038 & -0.0057 to +0.0129 & +0.0009 & -0.0016 to
+0.0038 \\
\end{longtable}
}

\textbf{Family evidence (C1).} Siblings from the same launch and breakup
that were catalogued, and had decayed, before the prediction time
improve TabM at 2016. At 2011 the point estimate is positive but the
interval includes zero.

\textbf{Physics summaries.} They help Kumo-Tabular at 2011, by a margin
that would not survive a multiplicity correction, and hurt TabM at 2016.
A plausible reason, not tested here: in 2016--2021 the physics expert's
own error grew with solar cycle 25, and a learner that leans on its
summaries inherits that error.

\textbf{Design properties.} They help at 2016.

\begin{figure}
\centering
\includegraphics[width=1\linewidth,height=\textheight,keepaspectratio,alt={Feature-group ablations of the selected learner (paired cluster bootstrap over launch families). In the first four rows a group is removed, so a positive change means the group helps. The last row instead adds the leaking 2026 radar cross-section, so there a positive change means the covariate makes the model worse.}]{artwork/Figure_9.png}
\caption{\textbf{Feature-group ablations of the selected learner}
(paired cluster bootstrap over launch families). In the first four rows
a group is removed, so a positive change means the group helps. The last
row instead \emph{adds} the leaking 2026 radar cross-section, so there a
positive change means the covariate makes the model
worse.}\label{fig:ablation}
\end{figure}

\subsection{Debris module: collision risk}\label{sec:res-cac}

\textbf{Setup.} The official Kelvins test set has 2,167 events, of which
150 are high-risk (Table \ref{tbl:cac}, Fig. \ref{fig:cac}).
Kumo-Tabular had the lowest out-of-fold pinball loss among the recent
learners and became the expert. The gate then put almost no weight on
the persistence expert (mean 0.0002), so the gated model is in effect
Kumo-Tabular.

\textbf{Probabilistic forecast.} Kumo-Tabular's final-risk distribution
has the lowest pinball loss. Paired bootstrap over test events (post
hoc; a negative difference favours Kumo-Tabular): against TabICLv2
-0.039 (-0.053 to -0.025), against XGBoost -0.202 (-0.233 to -0.171) and
against persistence -0.392 (-0.451 to -0.337).

{\def\LTcaptype{none} % do not increment counter
\begin{longtable}[]{@{}
  >{\raggedright\arraybackslash}p{(\linewidth - 6\tabcolsep) * \real{0.2500}}
  >{\raggedright\arraybackslash}p{(\linewidth - 6\tabcolsep) * \real{0.2500}}
  >{\raggedright\arraybackslash}p{(\linewidth - 6\tabcolsep) * \real{0.2500}}
  >{\raggedright\arraybackslash}p{(\linewidth - 6\tabcolsep) * \real{0.2500}}@{}}
\toprule\noalign{}
\begin{minipage}[b]{\linewidth}\raggedright
Model
\end{minipage} & \begin{minipage}[b]{\linewidth}\raggedright
Pinball loss
\end{minipage} & \begin{minipage}[b]{\linewidth}\raggedright
AUC-PR of P(high risk)
\end{minipage} & \begin{minipage}[b]{\linewidth}\raggedright
AUC-ROC of P(high risk)
\end{minipage} \\
\midrule\noalign{}
\endhead
\bottomrule\noalign{}
\endlastfoot
Kumo-Tabular & 0.653 & 0.733 & 0.970 \\
TabICLv2 & 0.692 & 0.665 & 0.961 \\
Best tree baseline (XGBoost) & 0.855 & 0.606 & 0.945 \\
Persistence & 1.045 & 0.664 & 0.944 \\
\end{longtable}
}

Its conformal 90 \% intervals cover 94.2 \% of all test events but only
48.0 \% of the high-risk ones. With 66 high-risk training events, the
group-conditional calibration is coarse exactly where it matters.

\textbf{Official point metric.} It does not reward these gains.

\begin{itemize}
\tightlist
\item
  Every model's median forecast falls below the high-risk threshold for
  most high-risk events, so its \(F_2\) is low: 129 of the 150 high-risk
  events receive a median below −6 from Kumo-Tabular, giving \(F_2\) =
  0.17. Most of those medians are not at the −30 floor but just below
  the threshold (41 lie between −10 and −6), so the floor is not the
  main cause (Fig. \ref{fig:cac-outputs}).
\item
  The latest-risk baseline scores \(L\) = 0.694, the published value.
\item
  Gating the latest-risk point with each model's calibrated P(high risk)
  brings the best models close to the baseline:

  \begin{itemize}
  \tightlist
  \item
    TabM: 0.698 (difference from the baseline, 95 \% CI -0.005 to
    +0.019);
  \item
    the gated model: 0.704 (-0.055 to +0.071).
  \end{itemize}
\item
  No model beats the baseline.
\end{itemize}

This agrees with the challenge, where few teams beat the latest risk,
and with a pre-registered re-evaluation
(\citeproc{ref-uriot2022spacecraft}{Uriot et al., 2022};
\citeproc{ref-vazquez2026latest}{Vazquez, 2026}).

\begin{longtable}[]{@{}
  >{\raggedright\arraybackslash}p{(\linewidth - 16\tabcolsep) * \real{0.1111}}
  >{\raggedright\arraybackslash}p{(\linewidth - 16\tabcolsep) * \real{0.1111}}
  >{\raggedright\arraybackslash}p{(\linewidth - 16\tabcolsep) * \real{0.1111}}
  >{\raggedright\arraybackslash}p{(\linewidth - 16\tabcolsep) * \real{0.1111}}
  >{\raggedright\arraybackslash}p{(\linewidth - 16\tabcolsep) * \real{0.1111}}
  >{\raggedright\arraybackslash}p{(\linewidth - 16\tabcolsep) * \real{0.1111}}
  >{\raggedright\arraybackslash}p{(\linewidth - 16\tabcolsep) * \real{0.1111}}
  >{\raggedright\arraybackslash}p{(\linewidth - 16\tabcolsep) * \real{0.1111}}
  >{\raggedright\arraybackslash}p{(\linewidth - 16\tabcolsep) * \real{0.1111}}@{}}
\caption{\textbf{ESA Kelvins Collision Avoidance Challenge, official
test set.} \(L = \text{MSE}_\text{HR}/F_2\) of the median (predictions
clipped at -6.001 as in the challenge); pinball loss of the final-risk
quantiles; AUC-PR and Brier score of P(final risk \textgreater= -6);
coverage of the conformal 90 \% interval over all events and over the
high-risk events.}\label{tbl:cac}\tabularnewline
\toprule\noalign{}
\begin{minipage}[b]{\linewidth}\raggedright
Model
\end{minipage} & \begin{minipage}[b]{\linewidth}\raggedright
L (clipped)
\end{minipage} & \begin{minipage}[b]{\linewidth}\raggedright
F2
\end{minipage} & \begin{minipage}[b]{\linewidth}\raggedright
MSE\_HR
\end{minipage} & \begin{minipage}[b]{\linewidth}\raggedright
Pinball
\end{minipage} & \begin{minipage}[b]{\linewidth}\raggedright
AUC-PR P(high)
\end{minipage} & \begin{minipage}[b]{\linewidth}\raggedright
Brier P(high)
\end{minipage} & \begin{minipage}[b]{\linewidth}\raggedright
Cov. 90 \%
\end{minipage} & \begin{minipage}[b]{\linewidth}\raggedright
Cov. 90 \%, high risk
\end{minipage} \\
\midrule\noalign{}
\endfirsthead
\toprule\noalign{}
\begin{minipage}[b]{\linewidth}\raggedright
Model
\end{minipage} & \begin{minipage}[b]{\linewidth}\raggedright
L (clipped)
\end{minipage} & \begin{minipage}[b]{\linewidth}\raggedright
F2
\end{minipage} & \begin{minipage}[b]{\linewidth}\raggedright
MSE\_HR
\end{minipage} & \begin{minipage}[b]{\linewidth}\raggedright
Pinball
\end{minipage} & \begin{minipage}[b]{\linewidth}\raggedright
AUC-PR P(high)
\end{minipage} & \begin{minipage}[b]{\linewidth}\raggedright
Brier P(high)
\end{minipage} & \begin{minipage}[b]{\linewidth}\raggedright
Cov. 90 \%
\end{minipage} & \begin{minipage}[b]{\linewidth}\raggedright
Cov. 90 \%, high risk
\end{minipage} \\
\midrule\noalign{}
\endhead
\bottomrule\noalign{}
\endlastfoot
Latest risk (challenge baseline) & 0.694 & 0.739 & 0.513 & --- & --- &
--- & --- & --- \\
Persistence + residual climatology & ∞ & 0.000 & 1.752 & 1.045 & 0.664 &
0.0523 & 0.910 & 0.593 \\
XGBoost (baseline) & ∞ & 0.000 & 1.752 & 0.855 & 0.606 & 0.0552 & 0.917
& 0.433 \\
LightGBM (baseline) & ∞ & 0.000 & 1.752 & 0.861 & 0.541 & 0.0601 & 0.905
& 0.253 \\
CatBoost (baseline) & ∞ & 0.000 & 1.752 & 0.869 & 0.542 & 0.0564 & 0.912
& 0.467 \\
Quantile MLP (baseline) & ∞ & 0.000 & 1.752 & 1.035 & 0.434 & 0.0653 &
0.936 & 0.853 \\
TabM (2025) & ∞ & 0.000 & 1.752 & 0.885 & 0.512 & 0.0591 & 0.919 &
0.727 \\
TabICLv2 (2026) & 30.059 & 0.057 & 1.728 & 0.692 & 0.665 & 0.0489 &
0.918 & 0.327 \\
Kumo-Tabular (2026) & 9.892 & 0.168 & 1.662 & 0.653 & 0.733 & 0.0453 &
0.942 & 0.480 \\
Persistence + expert, constant weight & 9.898 & 0.168 & 1.663 & 0.653 &
0.733 & 0.0453 & 0.935 & 0.480 \\
\textbf{Reliability-gated (COSMOS)} & 9.896 & 0.168 & 1.662 & 0.653 &
0.733 & 0.0453 & 0.940 & 0.480 \\
\end{longtable}

\begin{figure}
\centering
\includegraphics[width=1\linewidth,height=\textheight,keepaspectratio,alt={ESA Kelvins collision-risk test set. (a) Pinball loss of the final-risk distribution; (b) the official score of the latest-risk point forecast and of the same forecast gated by each model's calibrated P(high risk).}]{artwork/Figure_10.png}
\caption{\textbf{ESA Kelvins collision-risk test set.} (a) Pinball loss
of the final-risk distribution; (b) the official score of the
latest-risk point forecast and of the same forecast gated by each
model's calibrated P(high risk).}\label{fig:cac}
\end{figure}

\begin{figure}
\centering
\includegraphics[width=1\linewidth,height=\textheight,keepaspectratio,alt={What the collision-risk module outputs (official test set). (a) Final risk against the latest risk available at least two days before closest approach, the input of the persistence expert. (b) Final risk against Kumo-Tabular's median forecast. 129 of the 150 high-risk events (orange) receive a median below the −6 threshold, 41 of them between −10 and −6 and 37 within one decade of the −30 floor, so the median forecast reaches F\_2 = 0.17 on the official metric. Dashed lines mark the −6 threshold and the dotted line is equality. (c--e) Three high-risk test events drawn at random (fixed seed): the input messages, Kumo-Tabular's predictive quantiles (open circle: median) and the final risk, both drawn to the right of the input window.}]{artwork/Figure_11.png}
\caption{\textbf{What the collision-risk module outputs} (official test
set). (a) Final risk against the latest risk available at least two days
before closest approach, the input of the persistence expert. (b) Final
risk against Kumo-Tabular's median forecast. 129 of the 150 high-risk
events (orange) receive a median below the −6 threshold, 41 of them
between −10 and −6 and 37 within one decade of the −30 floor, so the
median forecast reaches \(F_2\) = 0.17 on the official metric. Dashed
lines mark the −6 threshold and the dotted line is equality. (c--e)
Three high-risk test events drawn at random (fixed seed): the input
messages, Kumo-Tabular's predictive quantiles (open circle: median) and
the final risk, both drawn to the right of the input
window.}\label{fig:cac-outputs}
\end{figure}

\subsection{Couplings}\label{sec:res-couplings}

\subsubsection{C2: from per-fragment re-entry to breakup-cloud
survival}\label{sec:res-c2}

\textbf{Setup.} C2 covers the test folds of both origins: 160
cloud-years in 27 breakup clouds, each with at least 30 catalogued
fragments. For each cloud we predict the number of its fragments still
in orbit one to five years after its first catalogued fragment, by
summing per-fragment survival probabilities, and compare it with the
catalogue count (Fig. \ref{fig:c2}).

\textbf{Results} (mean absolute error as a fraction of the cloud, with
95 \% intervals from a bootstrap over clouds; Supplementary Table S17):

{\def\LTcaptype{none} % do not increment counter
\begin{longtable}[]{@{}lllll@{}}
\toprule\noalign{}
Model & Raw reading & 95 \% CI & Calibrated reading & 95 \% CI \\
\midrule\noalign{}
\endhead
\bottomrule\noalign{}
\endlastfoot
Physics expert alone & 0.129 & 0.078--0.191 & 0.145 & 0.111--0.187 \\
TabICLv2 & 0.153 & 0.119--0.188 & 0.095 & 0.068--0.123 \\
Gated model & 0.155 & 0.121--0.192 & 0.122 & 0.085--0.161 \\
\end{longtable}
}

\begin{itemize}
\tightlist
\item
  \textbf{Raw reading.} The physics expert has the lowest point error.
  Its mean bias is +0.007, and 73.1 \% of cloud-years are within 10 \%.
  The learned experts under-predict survival on this reading, but the
  physics expert cannot be separated from the gated model (difference
  -0.025, 95 \% CI -0.084 to +0.033).
\item
  \textbf{Calibrated reading} (isotonic maps of \(P(T \le a)\) fitted on
  tune-B fragments). Calibration significantly lowers the error of
  TabICLv2 (+0.058, +0.034 to +0.082) and of the gated model (+0.033,
  +0.007 to +0.055). Calibrated TabICLv2 has the lowest point error,
  with 63.1 \% of cloud-years within 10 \%, but it cannot be separated
  from the raw physics expert (+0.035, -0.007 to +0.078). Calibration
  does not help the physics expert.
\end{itemize}

The calibrated reading and the cloud bootstrap were added after the raw
reading had been seen and criticised in an internal review; they are
post hoc.

\textbf{What C2 shows.} With 27 clouds, the errors of the best readings
are about a tenth of a cloud's size, and the uncertainty is too large to
rank the experts. How survival is read off the predictive distribution
matters for the learned experts. In the twin, cloud forecasts are sums
of the history-conditioned physics probabilities. T1 tests the
calibration of those probabilities on objects in orbit (Section
\ref{sec:res-t1}).

\begin{figure}
\centering
\includegraphics[width=1\linewidth,height=\textheight,keepaspectratio,alt={C2: predicted against observed fraction of each breakup cloud's fragments still in orbit, one point per cloud-year, (a) raw and (b) calibrated survival read-off.}]{artwork/Figure_12.png}
\caption{\textbf{C2: predicted against observed fraction of each breakup
cloud's fragments still in orbit}, one point per cloud-year, (a) raw and
(b) calibrated survival read-off.}\label{fig:c2}
\end{figure}

\subsubsection{C3: operational status as health evidence for the
re-entry regime}\label{sec:res-c3}

\textbf{Cohort.} On 2026-08-16, 13,794 Earth-orbiting payloads had
perigee below 600 km. In the following 72 days, 182 of them decayed; the
decays came from 86 of their 1,070 launches.

\textbf{Physics by status.} The drag physics over-predicts re-entry for
payloads reported operational:

{\def\LTcaptype{none} % do not increment counter
\begin{longtable}[]{@{}lllll@{}}
\toprule\noalign{}
Group & n & Decays & Mean predicted P & Observed rate \\
\midrule\noalign{}
\endhead
\bottomrule\noalign{}
\endlastfoot
Operational & 13,478 & 164 & 0.039 & 0.012 \\
Non-operational & 161 & 17 & 0.076 & 0.106 \\
\end{longtable}
}

\begin{itemize}
\tightlist
\item
  \textbf{Operational payloads.} The physics probabilities nevertheless
  rank them well (AUC-ROC 0.959), consistent with operators letting some
  low payloads decay while keeping others in orbit.
\item
  \textbf{Non-operational payloads.} The physics expert under-predicts
  somewhat, with AUC-ROC 0.935.
\end{itemize}

\textbf{As-of-date policies} (information at 2026-08-16 only; Fig.
\ref{fig:c3}):

{\def\LTcaptype{none} % do not increment counter
\begin{longtable}[]{@{}
  >{\raggedright\arraybackslash}p{(\linewidth - 4\tabcolsep) * \real{0.3333}}
  >{\raggedright\arraybackslash}p{(\linewidth - 4\tabcolsep) * \real{0.3333}}
  >{\raggedright\arraybackslash}p{(\linewidth - 4\tabcolsep) * \real{0.3333}}@{}}
\toprule\noalign{}
\begin{minipage}[b]{\linewidth}\raggedright
Policy
\end{minipage} & \begin{minipage}[b]{\linewidth}\raggedright
Brier score
\end{minipage} & \begin{minipage}[b]{\linewidth}\raggedright
ECE
\end{minipage} \\
\midrule\noalign{}
\endhead
\bottomrule\noalign{}
\endlastfoot
Drag physics for every payload & 0.0179 & 0.0267 \\
Previous-window base rate (0.0127) for every payload & 0.0130 &
0.0005 \\
Status-aware twin (base rate for operational payloads, physics
otherwise) & 0.0125 & 0.0008 \\
\end{longtable}
}

The status-aware twin has the lowest Brier score, but neither difference
is significant once launches are resampled:

{\def\LTcaptype{none} % do not increment counter
\begin{longtable}[]{@{}lll@{}}
\toprule\noalign{}
Status-aware twin minus & Brier difference & 95 \% CI \\
\midrule\noalign{}
\endhead
\bottomrule\noalign{}
\endlastfoot
Physics for every payload & -0.0054 & -0.0120 to 0.0006 \\
Base rate for every payload & -0.0005 & -0.0012 to 0.0001 \\
\end{longtable}
}

With 182 decays in one window, the test can show the direction of the
effect but cannot establish it. The window is also not fresh: an earlier
version of the project had fitted and tested re-entry classifiers on it
(Section \ref{sec:protocol}).

\textbf{In-window diagnostic (not a forecast).} We also fitted Platt
recalibrations on the follow-up window itself, cross-fitted by launch.
Recalibrating per status rather than for all payloads changes the Brier
score by -0.0004 (95 \% CI -0.0010 to 0.0001). Even with look-ahead,
then, status adds little once the physics probabilities are rescaled.
The regime split is what keeps the twin from issuing drag forecasts for
payloads whose operators control their orbits.

\begin{figure}
\centering
\includegraphics[width=1\linewidth,height=\textheight,keepaspectratio,alt={C3: mean predicted probability of re-entry within 72 days by operational status, against the observed rate (dashed). The in-window Platt map is fitted on the follow-up outcomes and is shown only as a reference.}]{artwork/Figure_13.png}
\caption{\textbf{C3: mean predicted probability of re-entry within 72
days} by operational status, against the observed rate (dashed). The
in-window Platt map is fitted on the follow-up outcomes and is shown
only as a reference.}\label{fig:c3}
\end{figure}

\subsubsection{T1: the twin's re-entry forecasts for objects in
orbit}\label{sec:res-t1}

\textbf{Cohort.} As of 2026-08-16, the twin held 13,186 fragments,
rocket stages and components with perigee below 2000 km. Within 72 days,
64 of them decayed, 29 within 30 days (Fig. \ref{fig:t1}).

\textbf{Results at 72 days} (expected re-entries, and their ratio to the
observed number):

{\def\LTcaptype{none} % do not increment counter
\begin{longtable}[]{@{}lllll@{}}
\toprule\noalign{}
Variant & Expected & Exp./obs. & 95 \% CI & Brier score \\
\midrule\noalign{}
\endhead
\bottomrule\noalign{}
\endlastfoot
History-conditioned physics & 87 & 1.37 & 1.20--1.63 & 0.00217 \\
Physics with class priors & 233 & 3.64 & 2.90--4.68 & 0.00625 \\
Lifetime pipeline (gated TabM) & 610 & 9.5 & 6.9--13.0 & 0.01344 \\
\end{longtable}
}

\begin{itemize}
\tightlist
\item
  \textbf{Lifetime pipeline.} Applied to current orbits, it
  over-predicts re-entry about ninefold, although it was the validated
  model at canonical-orbit epochs. Objects that have stayed in orbit for
  years are the survivors, with low area-to-mass ratios, and a model
  fitted at canonical-orbit epochs does not know this. Its TabM expert
  is refitted on the GPU, whose training is not bit-reproducible: an
  earlier run of the same configuration gave a ratio of 8.9.
\item
  \textbf{Physics with class priors.} It over-predicts about fourfold.
\item
  \textbf{History-conditioned physics.} Conditioning the area-to-mass
  ratio on the object's catalogue history cuts the excess to a factor of
  about 1.4. This is still significantly above one.
\end{itemize}

\textbf{Brier differences} (launch-cluster bootstrap):

{\def\LTcaptype{none} % do not increment counter
\begin{longtable}[]{@{}lll@{}}
\toprule\noalign{}
History-conditioned physics minus & Brier difference & 95 \% CI \\
\midrule\noalign{}
\endhead
\bottomrule\noalign{}
\endlastfoot
Physics with class priors & -0.00408 & -0.00598 to -0.00266 \\
Lifetime pipeline & -0.01127 & -0.01812 to -0.00632 \\
\end{longtable}
}

\textbf{Ranking.} Rank metrics use the raw probabilities. By AUC-PR the
history-conditioned variant ranks best at 72 days: 0.737, against 0.601
(class priors) and 0.616 (lifetime pipeline). By AUC-ROC the three are
within 0.01 of each other (0.991, 0.988 and 0.995), so all three
separate decaying objects from the rest about equally well. What differs
is mainly calibration. At 30 days the AUC-PR order is the same, and by
AUC-ROC the history-conditioned variant is slightly lower (0.982,
against 0.999 and 0.998).

\textbf{Payloads.}

\begin{itemize}
\tightlist
\item
  \emph{Operational payloads} (operator-controlled regime). The C3
  policy expected 170 decays against 164 observed. Drag physics would
  have expected 498. The base-rate policy cannot rank payloads, though:
  its AUC-ROC is 0.58, against 0.96 for drag physics.
\item
  \emph{Non-operational payloads.} The physics expert expected 11
  against 17 observed, the same under-prediction as in C3.
\end{itemize}

\textbf{The whole twin.} Over all 30,826 objects, the twin's per-regime
choice expected 272 re-entries against 246 observed (Brier 0.0065).

\textbf{Caveats.}

\begin{itemize}
\tightlist
\item
  T1 was specified before it was run, but it is not a blind test. An
  earlier version of the project had fitted and tested re-entry
  classifiers on the same window, and the window's aggregate decay count
  was seen while the lifetime pipeline's bias on current orbits was
  being diagnosed, before the history-conditioned variant was designed.
\item
  A defect in the payload prior was corrected after T1 had first been
  scored (Section \ref{sec:protocol}). Payload numbers changed; those of
  fragments, rocket stages and components did not.
\end{itemize}

\begin{figure}
\centering
\includegraphics[width=1\linewidth,height=\textheight,keepaspectratio,alt={T1: as-of-date test of the twin's re-entry forecasts for fragments, rocket stages and components in orbit on 2026-08-16. (a) Reliability at 72 days, in six fixed probability bins; marker area grows with the number of objects in the bin, bins with fewer than ten objects are dropped, and both axes are symmetric-log so that bins with no observed decay sit on the zero line. (b) Expected against observed re-entries.}]{artwork/Figure_14.png}
\caption{\textbf{T1: as-of-date test of the twin's re-entry forecasts}
for fragments, rocket stages and components in orbit on 2026-08-16. (a)
Reliability at 72 days, in six fixed probability bins; marker area grows
with the number of objects in the bin, bins with fewer than ten objects
are dropped, and both axes are symmetric-log so that bins with no
observed decay sit on the zero line. (b) Expected against observed
re-entries.}\label{fig:t1}
\end{figure}

\subsection{Health module}\label{sec:res-health-wrap}

\subsubsection{NASA SMAP/MSL telemetry and the companion ESA-ADB
study}\label{sec:res-health}

\textbf{Setup.} On the 81 channels (68 SMAP and 36 MSL anomaly
sequences), detectors are compared at a level and persistence fixed in
advance.

\textbf{Ranking} (Fig. \ref{fig:telemetry-benchmark}a). The published
telemanom scores (\citeproc{ref-hundman2018}{Hundman et al., 2018}) rank
anomalous steps best. Their per-channel mean AUC-PR is 0.576 (SMAP) and
0.473 (MSL). Used zero-shot, the Chronos-2 foundation model reaches
0.343 and 0.367. That is level with the ridge, PCA and nearest-neighbour
baselines, and above the random-score floor of 0.129 and 0.141.

\textbf{Alarms} (Fig. \ref{fig:telemetry-benchmark}b). On the ESA-ADB
corrected event-wise \(F_{0.5}\), which charges alarm time on nominal
data, the highest detector scores are PCA on SMAP (0.249) and the
nearest-neighbour detector on MSL (0.334). Chronos-2's residual alarms
are too frequent on MSL (508 false-positive segments), giving a
corrected \(F_{0.5}\) of 0.035 there, against 0.170 on SMAP. That
contrast is not a property of the spacecraft: one MSL channel
contributes 414 of those segments, and three contribute 95 \% of them.
They are channels whose training series is constant, so every
calibration score sits at the score floor and any non-zero test residual
raises an alarm.

\textbf{Why these alarm scores are weak evidence.}

\begin{itemize}
\tightlist
\item
  An always-on alarm reaches an uncorrected event \(F_1\) of 0.548
  (SMAP) and 0.537 (MSL), higher than every detector. Its corrected
  \(F_{0.5}\) is 0.
\item
  93.7 \% of labelled steps lie in the second half of their channel's
  test series. A label-free reference that keeps each detector's amount
  of alarm time but places it as one block at the end of the series
  scores higher on the corrected \(F_{0.5}\) than the detector itself in
  seven of the eight detector--spacecraft pairs; for example, 0.306
  against 0.249 for PCA on SMAP. The exception is the nearest-neighbour
  detector on MSL (0.334 against 0.281).
\end{itemize}

We therefore read skill against a null that keeps each detector's alarm
positions but reassigns them across channels. Every fitted detector
exceeds it on both spacecraft. For example, ridge's event \(F_1\) is
0.295 against 0.136 under the null (permutation \(p\) = 0.005). This
shows channel-specific detection, but a modest one: on this benchmark,
alarms that are merely late are rewarded about as much. Fig.
\ref{fig:smap-outputs} shows the scores and alarms on one typical
channel of each spacecraft.

\textbf{Companion ESA-ADB study.} The health module's alarm layer
(conformal p-values, Benjamini--Hochberg aggregation across channels and
persistence) was developed and tested in a separate, pre-registered
study on ESA-ADB (\citeproc{ref-kotowski2026esaadb}{Kotowski et al.,
2026}) by the same authors. On Mission 1, fitting an AR(8) forecaster
and selecting the alarm settings on the whole training half, rather than
on a short calibration window, raised the corrected event-wise
\(F_{0.5}\) from 0.074 to 0.428 (difference +0.353, 95 \% CI +0.206 to
+0.488), at about 1 false alarm per year. On Mission 2 the same contrast
reversed (-0.089, -0.143 to -0.049). Validation analyses of that study
attributed the Mission 1 gain to the labels used to select the alarm
settings, not to the fitted scores.

\textbf{Summary.} The recent forecasting foundation model does not
improve on simple detectors on SMAP/MSL, and the alarm protocol that the
twin's health module uses helped on one ESA-ADB mission and hurt on the
other. The health module is therefore the least validated of the three.

\begin{figure}
\centering
\includegraphics[width=1\linewidth,height=\textheight,keepaspectratio,alt={Telemetry benchmark on SMAP/MSL. (a) Per-channel mean AUC-PR of anomalous steps. (b) Corrected event-wise F0.5 of each detector's alarms (filled) and of a label-free reference that keeps the detector's amount of alarm time but places it as one block at the end of each series (open).}]{artwork/Figure_15.png}
\caption{\textbf{Telemetry benchmark on SMAP/MSL.} (a) Per-channel mean
AUC-PR of anomalous steps. (b) Corrected event-wise F0.5 of each
detector's alarms (filled) and of a label-free reference that keeps the
detector's amount of alarm time but places it as one block at the end of
each series (open).}\label{fig:telemetry-benchmark}
\end{figure}

\begin{figure}
\centering
\includegraphics[width=1\linewidth,height=\textheight,keepaspectratio,alt={What the health module outputs on one channel per spacecraft, chosen by rule as the median channel by Chronos-2 per-channel AUC-PR, taking the lower of the two middle channels where that number is even (SMAP E-6, MSL C-2). Top: test series and labelled anomalies. Middle: conformal p-values of PCA reconstruction and Chronos-2 against the channel's own training scores, with the alarm level. Bottom: labelled sequences and each detector's causal alarms. The training series of C-2 is constant, so every Chronos-2 calibration score sits at the score floor and any non-zero test residual alarms; Chronos-2 is in alarm from step 210 on, the most extreme Chronos-2 alarm behaviour of the 27 MSL channels. The recomputed alarms reproduce the benchmark's event counts for both channels.}]{artwork/Figure_16.png}
\caption{\textbf{What the health module outputs} on one channel per
spacecraft, chosen by rule as the median channel by Chronos-2
per-channel AUC-PR, taking the lower of the two middle channels where
that number is even (SMAP E-6, MSL C-2). Top: test series and labelled
anomalies. Middle: conformal p-values of PCA reconstruction and
Chronos-2 against the channel's own training scores, with the alarm
level. Bottom: labelled sequences and each detector's causal alarms. The
training series of C-2 is constant, so every Chronos-2 calibration score
sits at the score floor and any non-zero test residual alarms; Chronos-2
is in alarm from step 210 on, the most extreme Chronos-2 alarm behaviour
of the 27 MSL channels. The recomputed alarms reproduce the benchmark's
event counts for both channels.}\label{fig:smap-outputs}
\end{figure}

\subsection{The dated twin forecast of 2026-10-27}\label{sec:res-twin}

\textbf{Scope.} As of 2026-10-27 the twin holds 31,396 objects in orbit
with perigee below 2000 km:

\begin{itemize}
\tightlist
\item
  13,787 in the natural-decay regime;
\item
  16,386 operator-controlled payloads;
\item
  1,223 payloads without a reported status.
\end{itemize}

\textbf{History conditioning} (Fig. \ref{fig:twin-state} shows three
examples). Among fragments, rocket stages and components:

\begin{itemize}
\tightlist
\item
  for 12,606, the catalogue history determined the area-to-mass
  posterior (effective sample size of at least 4);
\item
  for 410, survival alone was used;
\item
  for 56, no prior sample survived, and the longest survivors were used;
\item
  166 objects without a usable canonical orbit kept the class prior.
\end{itemize}

\textbf{Expected re-entries in the next 30 days} (natural-decay regime):
24.5. For fragments, rocket stages and components, the three variants
differ by an order of magnitude:

{\def\LTcaptype{none} % do not increment counter
\begin{longtable}[]{@{}ll@{}}
\toprule\noalign{}
Variant & Expected re-entries, 30 days \\
\midrule\noalign{}
\endhead
\bottomrule\noalign{}
\endlastfoot
History-conditioned physics & 21.6 \\
Physics with class priors & 76.3 \\
Gated lifetime pipeline, fitted at canonical-orbit epochs & 229.1 \\
\end{longtable}
}

T1 indicates which of these to believe (Section \ref{sec:res-t1}). Fig.
\ref{fig:twin-forecast} maps the forecast and lists the natural-decay
objects with the highest 30-day probability.

\textbf{Clouds and verification.} The forecast includes breakup-cloud
forecasts for the 61 clouds with at least 30 fragments in orbit. Its
record file is fixed by its SHA-256 hash,
\texttt{19a1fede1fd4871f380ff802f285d2ad03732f6cd3179790e67336c7d202a6b5},
so that later catalogue snapshots can score it without the possibility
of revision. An earlier forecast of the same date (hash beginning
\texttt{fad1c619}) contained the payload-prior defect described in
Section \ref{sec:protocol}. It was replaced before any later snapshot
existed, and it is archived with the code.

\begin{figure}
\centering
\includegraphics[width=1\linewidth,height=\textheight,keepaspectratio,alt={Twin state estimation, for three objects drawn at random (one per class; fixed seed) among natural-decay objects whose catalogue history determined the posterior (effective sample size at least 4) and whose history-conditioned median forecast is within five years. (a) Prior area-to-mass samples (grey: the object's design value where GCAT has one, otherwise a draw from its class prior) and the posterior after conditioning on the change from the canonical to the current orbit (orange). (b) The re-entry forecasts the same samples produce. The recomputed posteriors reproduce the stored forecasts.}]{artwork/Figure_17.png}
\caption{\textbf{Twin state estimation}, for three objects drawn at
random (one per class; fixed seed) among natural-decay objects whose
catalogue history determined the posterior (effective sample size at
least 4) and whose history-conditioned median forecast is within five
years. (a) Prior area-to-mass samples (grey: the object's design value
where GCAT has one, otherwise a draw from its class prior) and the
posterior after conditioning on the change from the canonical to the
current orbit (orange). (b) The re-entry forecasts the same samples
produce. The recomputed posteriors reproduce the stored
forecasts.}\label{fig:twin-state}
\end{figure}

\begin{figure}
\centering
\includegraphics[width=1\linewidth,height=\textheight,keepaspectratio,alt={The dated twin forecast of 2026-10-27. (a) Current orbits of the objects in the twin, by regime; the regimes overlap, and the largest is drawn first. (b) The twin's probability of re-entry within 30 days against current perigee, natural-decay regime. (c) The 15 natural-decay objects with the highest probability, ranked by that probability and then by the median forecast: median and 90 \% Monte-Carlo interval of the forecast re-entry time (history-conditioned physics for fragments, rocket stages and components; class-prior physics for non-operational payloads). Two rows have a degenerate interval at exactly 10 days: their perigee is already within 1 km of the 120 km re-entry threshold, where the integrator's maximum step of 10 days sets the resolution.}]{artwork/Figure_18.png}
\caption{\textbf{The dated twin forecast of 2026-10-27.} (a) Current
orbits of the objects in the twin, by regime; the regimes overlap, and
the largest is drawn first. (b) The twin's probability of re-entry
within 30 days against current perigee, natural-decay regime. (c) The 15
natural-decay objects with the highest probability, ranked by that
probability and then by the median forecast: median and 90 \%
Monte-Carlo interval of the forecast re-entry time (history-conditioned
physics for fragments, rocket stages and components; class-prior physics
for non-operational payloads). Two rows have a degenerate interval at
exactly 10 days: their perigee is already within 1 km of the 120 km
re-entry threshold, where the integrator's maximum step of 10 days sets
the resolution.}\label{fig:twin-forecast}
\end{figure}

\subsubsection{Computational cost}\label{sec:res-cost}

Times are measured on the origin-2016 run (25,334 training objects,
2,993 test objects; Fig. \ref{fig:cost}). Hardware:

\begin{itemize}
\tightlist
\item
  CPU learners, TabICLv2 and the physics expert: a two-socket Intel Xeon
  6980P node (512 hardware threads).
\item
  TabM and Kumo-Tabular: an Intel Arc Pro B60 GPU.
\end{itemize}

{\def\LTcaptype{none} % do not increment counter
\begin{longtable}[]{@{}
  >{\raggedright\arraybackslash}p{(\linewidth - 4\tabcolsep) * \real{0.3333}}
  >{\raggedright\arraybackslash}p{(\linewidth - 4\tabcolsep) * \real{0.3333}}
  >{\raggedright\arraybackslash}p{(\linewidth - 4\tabcolsep) * \real{0.3333}}@{}}
\toprule\noalign{}
\begin{minipage}[b]{\linewidth}\raggedright
Model
\end{minipage} & \begin{minipage}[b]{\linewidth}\raggedright
Fit and predict per seed (s)
\end{minipage} & \begin{minipage}[b]{\linewidth}\raggedright
One call over the test cohort (s)
\end{minipage} \\
\midrule\noalign{}
\endhead
\bottomrule\noalign{}
\endlastfoot
Physics expert (32 Monte-Carlo samples) & 128 & 34 \\
XGBoost & 26 & 0.029 \\
TabM & 60 & 0.063 \\
Kumo-Tabular & 94 & 95 \\
TabICLv2 & 915 & 617 \\
\end{longtable}
}

\textbf{What is timed.} The first column is the wall-clock of the
fit-and-predict loop divided by the five seeds: it contains four refits
(the seed-0 fit happens during tuning and is not timed) and the
predictions of all five seeds on the tuning and test cohorts. It is
therefore an upper bound on fitting alone, and for the in-context models
it is dominated by prediction. The second column is one call predicting
the test cohort.

\textbf{In-context foundation models.} Their prediction cost includes
encoding the training context, up to 16,000 rows for Kumo-Tabular and
25,334 for TabICLv2 on CPU. Predicting the same cohort takes them three
to five orders of magnitude longer than a tree ensemble (Kumo-Tabular 95
s and TabICLv2 617 s, against 0.029 s for XGBoost and 0.005 s for
CatBoost). Because the context encoding is a fixed cost per call rather
than per object, these times cannot be rescaled to a different number of
objects, and we do not extrapolate them to the full catalogue; both
models predicted the whole test cohort in one call.

\textbf{The physics features are part of every learner's cost.} All
learners take five summaries of the physics expert's Monte-Carlo
distribution as inputs, so a deployed pipeline also pays that run: 34 s
for the 2,993 test objects, which dominates the tree ensembles' own
prediction time by three orders of magnitude. The second column above
times the learner alone.

\begin{figure}
\centering
\includegraphics[width=1\linewidth,height=\textheight,keepaspectratio,alt={Computational cost, origin 2016, on log scales. (a) Wall-clock of the fit-and-predict loop per seed: four refits (the seed-0 fit happens during tuning and is not timed) plus the predictions of all five seeds on the tuning and test cohorts, so it bounds fitting alone from above and for the in-context models is dominated by prediction; for the physics expert, its Monte-Carlo run over the training set. (b) One call predicting the test cohort, which for the in-context models includes encoding the training context, a fixed cost per call rather than per object.}]{artwork/Figure_19.png}
\caption{\textbf{Computational cost}, origin 2016, on log scales. (a)
Wall-clock of the fit-and-predict loop per seed: four refits (the seed-0
fit happens during tuning and is not timed) plus the predictions of all
five seeds on the tuning and test cohorts, so it bounds fitting alone
from above and for the in-context models is dominated by prediction; for
the physics expert, its Monte-Carlo run over the training set. (b) One
call predicting the test cohort, which for the in-context models
includes encoding the training context, a fixed cost per call rather
than per object.}\label{fig:cost}
\end{figure}

\section{Discussion}\label{sec:discussion}

\textbf{Which re-entry expert is best depends on the test period.}

\begin{itemize}
\tightlist
\item
  For 2016--2021, two in-context tabular foundation models released in
  2026, TabICLv2 and Kumo-Tabular, were significantly more accurate than
  the physics expert, the gated model, the selected TabM and the best
  tree ensemble. They also had the best concordance and near-nominal raw
  coverage.
\item
  For 2011--2015, the physics expert alone had the lowest point
  estimate, but it could not be separated from the foundation models or
  from the gated model. It had the lowest loss for fragments in that
  period. The one significant difference among them is that
  Kumo-Tabular, the learner selected there, is worse than the gated
  model that fuses it.
\end{itemize}

TabICLv2 was thus best or tied in both periods. One change between the
periods is measured: replaying the realised F10.7 lowers the physics
expert's error much more in the rising phase of cycle 25 than in cycle
24. Other changes in the test cohorts were not isolated here. Two test
periods, of which several comparisons were added after the first results
were seen, are not enough to say whether the foundation models will keep
their lead.

\textbf{Selecting a learner on a tuning cohort five to seven years older
than the test cohort did not reliably pick the model that generalised
best.}

\begin{itemize}
\tightlist
\item
  At 2016, TabM led on tune-A by a wide margin and was significantly
  worse than both foundation models on test.
\item
  At 2011, the small tune-A set (248 objects) chose Kumo-Tabular, which
  is tied with TabICLv2 on test.
\item
  The first learner set, recorded as a deviation, selected its MLP at
  both origins, and the MLP and TabM failed on 2011 fragments.
\end{itemize}

The temporal protocol is what exposes this. A random split would have
hidden it.

\textbf{The reliability gate helped when the physics expert was
reliable, and it did not beat the best single expert.}

\begin{itemize}
\tightlist
\item
  At 2011 it significantly improved every learner to which it gave
  non-negligible weight (post-hoc sensitivity analysis). The reported
  2011 gated model depends on the change of the candidate pool: with the
  first learner set, its gated model was much worse.
\item
  At 2016, when solar activity degraded the physics expert, it improved
  none and cost a little. The cost was not statistically significant.
\end{itemize}

The gate is fitted on out-of-fold predictions of the training period, so
it cannot anticipate that a learner will degrade on a future cohort. Its
descriptors (area-to-mass source, the physics expert's spread, orbit,
class) describe the physics expert's reliability, not the learner's
reliability under shift. A gate that also sees descriptors of covariate
shift is a natural next step, but it would have to be validated on a
third test period.

\textbf{The radar-cross-section leak is a warning for catalogue-based
learning.}

\begin{itemize}
\tightlist
\item
  A covariate that ``is a size measurement'' encodes the outcome through
  its availability (adjusted odds ratio 3.6, cluster-robust 95 \% CI
  2.1--6.2).
\item
  Adding it changed the test pinball loss by amounts within noise at
  both origins, so point accuracy would not have revealed the leak.
\item
  At 2016, where its missingness pattern differs from training, it broke
  interval coverage.
\end{itemize}

Any study that joins catalogue snapshots taken long after the prediction
time should test the missingness of each covariate against the outcome.
Our code includes such a check.

\textbf{What the couplings show.}

\begin{itemize}
\tightlist
\item
  \textbf{Family evidence (C1)} improved the learner in one of two
  periods.
\item
  \textbf{Cloud survival (C2)} is accurate to about a tenth of a cloud's
  size for the best readings, but with 27 clouds the experts cannot be
  ranked.
\item
  \textbf{Operational status (C3)} marks a regime in which drag physics
  over-predicts re-entry by a factor of about three. A status-aware
  policy lowered the Brier score in one as-of-date window, but not
  significantly once outcomes are clustered by launch.
\item
  \textbf{The twin's current-state forecasts (T1)} are where the
  coupling of history and physics matters most. Conditioning on the
  catalogue history reduced the over-prediction of 72-day re-entries
  from 9.5-fold (lifetime pipeline) and 3.6-fold (class priors) to
  1.37-fold, in one window that had been analysed earlier in the
  project.
\end{itemize}

C1, C2 and C3 are small or uncertain effects. T1 is a large effect, but
it comes from one window that had been analysed before. Together they
are what a digital twin that only joins separate tools cannot provide: a
measured account of which links between the modules carry information.

\subsection{Limitations and threats to validity}\label{sec:limits}

\textbf{One orbit per object.}

\begin{itemize}
\tightlist
\item
  Orbital-element histories at catalogue scale are not legitimately
  public: Space-Track requires an account and forbids redistribution.
\item
  The re-entry module therefore predicts from the canonical orbit, and
  the twin uses two orbits: canonical and current.
\item
  Operational tools that fit the ballistic coefficient to weeks of TLE
  data (\citeproc{ref-gondelach2017ballistic}{Gondelach et al., 2017};
  \citeproc{ref-pardini2008impact}{Pardini and Anselmo, 2008}) use far
  more information near re-entry. Our months-to-years forecasts are not
  comparable with their final-days accuracy.
\end{itemize}

\textbf{Validation of current-state forecasts.}

\begin{itemize}
\tightlist
\item
  The lifetime study validates forecasts at canonical-orbit epochs.
\item
  Forecasts for objects already in orbit are validated only by T1. T1
  covers one 72-day window, and it is not blind: an earlier version of
  the project fitted and tested re-entry classifiers on the same window,
  and its aggregate decay count was seen before the history-conditioned
  variant was designed.
\item
  The hashed forecast of 2026-10-27 is a blind test that a later
  catalogue snapshot will score.
\end{itemize}

\textbf{Priors built from catalogue values.}

\begin{itemize}
\tightlist
\item
  The class prior of characteristic length comes from training objects
  that received a radar-cross-section value, so it is biased towards
  long-lived objects.
\item
  Among components, the availability of design data is associated with
  early decay (adjusted odds ratio 3.4). We treat it as known at launch,
  because GCAT design values come from launch documentation. The
  \texttt{no\_design} ablation bounds its influence.
\end{itemize}

\textbf{Resolution of the decay integrator.} A step may not drop the
perigee by more than 2 \% of its height above 120 km, with a floor of 1
km, and may not exceed 10 days. For an object whose perigee is already
within 1 km of the 120 km threshold, the first step can therefore be a
full 10 days, and the forecast is reported at the end of that step. Two
of the 13,787 natural-decay objects of the dated forecast have a
degenerate 90 \% interval at exactly 10 days for this reason. The
affected objects are within days of re-entry in any case, and no
reported comparison depends on them.

\textbf{Physics simplifications.}

\begin{itemize}
\tightlist
\item
  The density is a global NRLMSIS 2.1 average without latitude,
  local-time or seasonal variation.
\item
  The atmosphere is exponential about perigee within an orbit.
\item
  Lunisolar and J2 effects on perigee are not modelled. They matter for
  transfer-orbit stages.
\item
  Solar activity is held at its value at the prediction time. The oracle
  diagnostic shows how much of the error this causes (Section
  \ref{sec:results}).
\end{itemize}

\textbf{Small calibration sets at the 2011 origin.} The 2004--2005
tuning cohort is small, so model selection and calibration at that
origin are noisy.

\textbf{Test-set reuse and post-hoc analyses.}

\begin{itemize}
\tightlist
\item
  The re-entry and collision-risk test sets were scored more than once
  during development: before an internal review, with a first learner
  set, and in smoke tests of the new learners.
\item
  The pool of eligible learners was changed after the first set's test
  results were known.
\item
  The head-to-head comparisons, the gate applied to every learner and
  the calibrated C2 reading were added after results were seen. Their
  intervals are nominal and uncorrected for multiplicity.
\end{itemize}

All of this is recorded in Section \ref{sec:protocol} and in the
released logs. It weakens the confirmatory value of the comparisons, and
a fresh test period is needed to confirm them.

\textbf{Bit-level reproducibility.} Seeded training of TabM on the GPU
is not bit-reproducible; on CPU it is. The re-entry study reuses cached
predictions, so its numbers are fixed, but a refit gives slightly
different results: two runs of the T1 lifetime-pipeline variant gave
expected/observed ratios of 8.9 and 9.5. Chronos-2 inference on the GPU
varies at the same level; repeating the SMAP/MSL run moved its
per-channel AUC-PR by at most \(2 \times 10^{-5}\) and left every
reported value unchanged at three decimals.

\textbf{Health evidence is a proxy.}

\begin{itemize}
\tightlist
\item
  Catalogue operational status is the only health evidence available for
  every payload. It is not a measurement of manoeuvre capability.
\item
  No public dataset links telemetry, conjunction messages and catalogue
  objects, so the health-to-debris coupling (C4) is a demonstration
  only.
\end{itemize}

\textbf{Event metrics for telemetry.} Event-wise metrics can be
satisfied by trivial alarm rules, so they are reported next to trivial
references and with the corrected precision of ESA-ADB. Even so, a
reference that only places alarms late in each series beats most
detectors on SMAP/MSL. The alarm protocol of the health module helped on
one ESA-ADB mission and hurt on the other, so the health module is the
least validated part of the twin.

\section{Conclusions}\label{sec:conclusions}

COSMOS-AI joins satellite health monitoring, orbital debris analysis and
atmospheric re-entry prediction in one digital twin. Each module is
validated on public data, and the couplings are tested wherever a real
test exists. Five findings stand out.

\begin{enumerate}
\def\labelenumi{\arabic{enumi}.}
\tightlist
\item
  \textbf{Re-entry prediction across solar-cycle phases.}

  \begin{itemize}
  \tightlist
  \item
    Two 2026 tabular foundation models were significantly the most
    accurate in 2016--2021.
  \item
    In 2011--2015 TabICLv2 was statistically tied with the physics
    expert and the gated model; Kumo-Tabular was significantly worse
    than the gated model that fuses it.
  \item
    Selection on an older tuning cohort chose a significantly worse
    learner in 2016.
  \end{itemize}

  The reliability gate improved the learners it weighted at the 2011
  origin and was neutral to slightly harmful at 2016. It did not beat
  the best single expert.
\item
  \textbf{Catalogue covariates can carry the outcome in their
  availability.} Whether an object has a radar cross-section is strongly
  associated with early decay. Using it broke interval coverage in the
  period whose missingness pattern differed from training, while its
  effect on point accuracy was within noise.
\item
  \textbf{For objects already in orbit, the twin's state estimate
  matters more than the choice of learner.} Conditioning each object's
  area-to-mass ratio on its catalogue history reduced the
  over-prediction of 72-day re-entries from 9.5-fold (the lifetime
  pipeline) and 3.6-fold (physics with class priors) to 1.4-fold.
\item
  \textbf{Operational status separates payloads whose re-entry drag
  physics cannot predict.} Using it lowered the Brier score in one
  72-day window, but not significantly.
\item
  \textbf{Conjunction and telemetry: the newest models did not beat the
  established baselines on the operational metrics.}

  \begin{itemize}
  \tightlist
  \item
    On the ESA collision-risk benchmark, a 2026 tabular foundation model
    gave significantly the best probabilistic forecasts. Neither it nor
    any other model beat the latest-risk forecast on the official point
    metric.
  \item
    A zero-shot time-series foundation model did not beat simple
    detectors on SMAP/MSL.
  \end{itemize}
\end{enumerate}

\textbf{Limits.}

\begin{itemize}
\tightlist
\item
  The twin rests on single catalogue orbits rather than orbital
  histories.
\item
  Findings 3 and 4 come from one as-of-date window that had been
  analysed earlier in the project, so they are not blind tests.
\item
  The test sets were scored more than once during development, and
  several analyses were added after results were seen (Section
  \ref{sec:protocol}).
\end{itemize}

\textbf{Verification.} The hashed forecast of 2026-10-27 is the first
blind test of the twin and will be scored against later catalogue
snapshots. The code and derived data needed to repeat every number are
released with the paper.

\section*{Declarations}\label{declarations}
\addcontentsline{toc}{section}{Declarations}

\textbf{CRediT authorship contribution statement.} {[}Author 1{]}:
Conceptualization, Methodology, Software, Validation, Formal analysis,
Investigation, Data curation, Writing -- original draft, Writing --
review \& editing, Visualization. {[}Author 2{]}: Supervision, Writing
-- review \& editing. {[}To be completed by the authors.{]}

\textbf{Data availability.} All data used are public. GCAT
(\citeproc{ref-mcdowell2020gcat}{McDowell, 2020}) is available under
CC-BY-4.0; CelesTrak catalogue and space-weather files
(\citeproc{ref-celestrak2026satcat}{Kelso, 2026}) under CelesTrak's
usage policy; the ESA Kelvins Collision Avoidance Challenge data
(\citeproc{ref-uriot2021cacdataset}{Uriot et al., 2021}) under
CC-BY-4.0; ESA-ADB (\citeproc{ref-esaadb_dataset}{De Canio et al.,
2024}) (companion study) and NASA SMAP/MSL
(\citeproc{ref-hundman2018}{Hundman et al., 2018}) under the terms
stated by their authors. The derived cohort table, the NRLMSIS density
table, all predictions and result files, and the dated twin forecast
will be deposited in Zenodo with a DOI before publication {[}DOI to be
added{]}. No Space-Track data are used or redistributed. Pretrained
model weights are obtained from their publishers under their own
licences:

\begin{itemize}
\tightlist
\item
  TabICLv2: BSD-3-Clause;
\item
  Kumo-Tabular: OpenMDW-1.1;
\item
  Chronos-2: Apache-2.0;
\item
  NRLMSIS 2.1, via pymsis.
\end{itemize}

\textbf{Code availability.} The code (data fetchers with checksummed
manifests, the re-entry, conjunction and telemetry modules, the twin
engine, and the scripts that produce every table and figure from the
result files) is available on GitHub
(\citeproc{ref-cosmos_software}{{[}Author name(s){]}, 2026}); an
archived release with a DOI will be created before publication.

\textbf{Funding.} {[}To be completed by the authors.{]}

\textbf{Declaration of competing interest.} {[}To be completed by the
authors.{]}

\textbf{Declaration of generative AI and AI-assisted technologies in the
manuscript preparation process.} During the preparation of this work the
author(s) used Claude (Anthropic) to assist with writing analysis and
plotting code and with drafting and editing the text. After using this
tool, the author(s) reviewed and edited the content as needed and
take(s) full responsibility for the content of the published article.
All figures were rendered with Matplotlib from the study's result files;
no figure was generated or altered by a generative AI model. {[}The
authors must review this statement and adjust it to their actual use.{]}

\textbf{What was not validated.} The C4 coupling (health × debris →
collision response) is a demonstration: no public dataset pairs
conjunction messages with the health of the same spacecraft, and
operational status is a proxy for manoeuvre capability, not a
measurement of it. The twin forecast as of 2026-10-27 is a prospective
forecast; it will be verified against later catalogue snapshots. The
as-of-date test T1 covers one 72-day window that had been analysed
earlier in the project, so it is not a blind test (Section
\ref{sec:protocol}).

\section*{References}\label{references}
\addcontentsline{toc}{section}{References}

\protect\phantomsection\label{refs}
\begin{CSLReferences}{1}{0}
\bibitem[\citeproctext]{ref-angelopoulos2023gentle}
Angelopoulos, A.N., Bates, S., 2023. Conformal prediction: A gentle
introduction. Foundations and Trends in Machine Learning 16, 494--591.
\url{https://doi.org/10.1561/2200000101}

\bibitem[\citeproctext]{ref-ansari2025chronos2}
Ansari, A.F., Shchur, O., Küken, J., Auer, A., Han, B., Mercado, P.,
Rangapuram, S.S., Shen, H., Stella, L., Zhang, X., Goswami, M., Kapoor,
S., Maddix, D.C., Guerron, P., Hu, T., Yin, J., Erickson, N., Desai,
P.M., Wang, H., Rangwala, H., Karypis, G., Wang, Y., Bohlke-Schneider,
M., 2025. {Chronos-2}: From univariate to universal forecasting.
\url{https://arxiv.org/abs/2510.15821}

\bibitem[\citeproctext]{ref-cosmos_software}
{[}software{]} {[}Author name(s){]}, 2026. {COSMOS-AI: spacecraft
telemetry alarm-protocol study code}. GitHub, version commit 2525367.
\url{https://github.com/AnuSajikumar6264/Cosmos-AI-}

\bibitem[\citeproctext]{ref-celletti2021reconnecting}
Celletti, A., Pucacco, G., Vartolomei, T., 2021. Reconnecting groups of
space debris to their parent body through proper elements. Scientific
Reports 11, 22676. \url{https://doi.org/10.1038/s41598-021-02010-x}

\bibitem[\citeproctext]{ref-chen2016xgboost}
Chen, T., Guestrin, C., 2016. {XGBoost}: A scalable tree boosting
system, in: Proceedings of the 22nd {ACM SIGKDD} International
Conference on Knowledge Discovery and Data Mining. ACM, pp. 785--794.
\url{https://doi.org/10.1145/2939672.2939785}

\bibitem[\citeproctext]{ref-esaadb_dataset}
{[}dataset{]} De Canio, G., Kotowski, K., Haskamp, C., 2024. {ESA
Anomaly Dataset}. Zenodo, version 1.0.
\url{https://doi.org/10.5281/zenodo.12528696}

\bibitem[\citeproctext]{ref-efthymiopoulos2021origin}
Efthymiopoulos, C., Märtens, M., Gaitanas, M., Izzo, D., 2021. Space
debris: The origin dataset. \url{https://doi.org/10.5281/zenodo.4986292}

\bibitem[\citeproctext]{ref-emmert2021nrlmsis}
Emmert, J.T., Drob, D.P., Picone, J.M., Siskind, D.E., Jones, M.,
Mlynczak, M.G., Bernath, P.F., Chu, X., Doornbos, E., Funke, B.,
Goncharenko, L.P., Hervig, M.E., Schwartz, M.J., Sheese, P.E., Vargas,
F., Williams, B.P., Yuan, T., 2021. {NRLMSIS} 2.0: A whole-atmosphere
empirical model of temperature and neutral species densities. Earth
Space Sci. 8, e2020EA001321. \url{https://doi.org/10.1029/2020EA001321}

\bibitem[\citeproctext]{ref-frey2019extension}
Frey, S., Colombo, C., Lemmens, S., 2019. Extension of the king-hele
orbit contraction method for accurate, semi-analytical propagation of
non-circular orbits. Advances in Space Research 64, 1--17.
\url{https://doi.org/10.1016/j.asr.2019.03.016}

\bibitem[\citeproctext]{ref-gondelach2017ballistic}
Gondelach, D.J., Armellin, R., Lidtke, A.A., 2017. Ballistic coefficient
estimation for reentry prediction of rocket bodies in eccentric orbits
based on TLE data. Mathematical Problems in Engineering 2017, 7309637.
\url{https://doi.org/10.1155/2017/7309637}

\bibitem[\citeproctext]{ref-gorishniy2025tabm}
Gorishniy, Y., Kotelnikov, A., Babenko, A., 2025. {TabM}: Advancing
tabular deep learning with parameter-efficient ensembling, in:
International Conference on Learning Representations (ICLR).
\url{https://arxiv.org/abs/2410.24210}

\bibitem[\citeproctext]{ref-harrell1982evaluating}
Harrell, F.E., Califf, R.M., Pryor, D.B., Lee, K.L., Rosati, R.A., 1982.
Evaluating the yield of medical tests. JAMA 247, 2543--2546.
\url{https://doi.org/10.1001/jama.1982.03320430047030}

\bibitem[\citeproctext]{ref-hundman2018}
Hundman, K., Constantinou, V., Laporte, C., Colwell, I., Soderstrom, T.,
2018. Detecting spacecraft anomalies using {LSTMs} and nonparametric
dynamic thresholding, in: Proceedings of the 24th {ACM SIGKDD}
International Conference on Knowledge Discovery \& Data Mining. ACM, pp.
387--395. \url{https://doi.org/10.1145/3219819.3219845}

\bibitem[\citeproctext]{ref-johnson2001nasa}
Johnson, N.L., Krisko, P.H., Liou, J.-C., Anz-Meador, P.D., 2001.
{NASA}'s new breakup model of {EVOLVE} 4.0. Advances in Space Research
28, 1377--1384. \url{https://doi.org/10.1016/S0273-1177(01)00423-9}

\bibitem[\citeproctext]{ref-jung2021recurrent}
Jung, O., Seong, J., Jung, Y., Bang, H., 2021. Recurrent neural network
model to predict re-entry trajectories of uncontrolled space objects.
Advances in Space Research 68, 2515--2529.
\url{https://doi.org/10.1016/j.asr.2021.04.041}

\bibitem[\citeproctext]{ref-ke2017lightgbm}
Ke, G., Meng, Q., Finley, T., Wang, T., Chen, W., Ma, W., Ye, Q., Liu,
T.-Y., 2017. {LightGBM}: A highly efficient gradient boosting decision
tree, in: Advances in Neural Information Processing Systems 30.
\url{https://proceedings.neurips.cc/paper/2017/hash/6449f44a102fde848669bdd9eb6b76fa-Abstract.html}

\bibitem[\citeproctext]{ref-celestrak2026satcat}
Kelso, T.S., 2026. {SATCAT} format documentation and {CelesTrak} usage
policy.

\bibitem[\citeproctext]{ref-kim2022rigorous}
Kim, S., Choi, K., Choi, H.-S., Lee, B., Yoon, S., 2022. Towards a
rigorous evaluation of time-series anomaly detection. Proc. {AAAI} Conf.
Artif. Intell. 36, 7194--7201.
\url{https://doi.org/10.1609/aaai.v36i7.20680}

\bibitem[\citeproctext]{ref-kinghele1987satellite}
King-Hele, D., 1987. Satellite orbits in an atmosphere: Theory and
applications. Blackie.

\bibitem[\citeproctext]{ref-kotowski2026esaadb}
Kotowski, K., Haskamp, C., Andrzejewski, J., Ruszczak, B., Nalepa, J.,
Lakey, D., Collins, P., Martínez-Heras, J., De Canio, G., 2026.
{European Space Agency} dataset and benchmark for real-world anomaly
detection in spacecraft time series. Journal of Data-centric Machine
Learning Research 3, 1--76.
\url{https://data.mlr.press/assets/pdf/v03-23.pdf}

\bibitem[\citeproctext]{ref-licata2022uncertainty}
Licata, R.J., Mehta, P.M., 2022. Uncertainty quantification techniques
for data-driven space weather modeling: Thermospheric density
application. Scientific Reports 12, 7256.
\url{https://doi.org/10.1038/s41598-022-11049-3}

\bibitem[\citeproctext]{ref-lichtendahl2013average}
Lichtendahl, K.C., Grushka-Cockayne, Y., Winkler, R.L., 2013. Is it
better to average probabilities or quantiles? Manag. Sci. 59,
1594--1611. \url{https://doi.org/10.1287/mnsc.1120.1667}

\bibitem[\citeproctext]{ref-liu2024digital}
Liu, W., Wu, M., Wan, G., Xu, M., 2024. Digital twin of space
environment: Development, challenges, applications, and future outlook.
Remote Sensing 16, 3023. \url{https://doi.org/10.3390/rs16163023}

\bibitem[\citeproctext]{ref-lotti2025exploring}
Lotti, A., Modenini, D., Tortora, P., 2025. Exploring physics-informed
neural networks for orbital lifetime estimation, in: 11th European
Conference for Aeronautics and Space Sciences (EUCASS 2025). EUCASS.
\url{https://doi.org/10.13009/EUCASS2025-698}

\bibitem[\citeproctext]{ref-mcdowell2020gcat}
McDowell, J.C., 2020. General catalog of artificial space objects.
\url{https://planet4589.org/space/gcat}

\bibitem[\citeproctext]{ref-nvidia2026kumo}
NVIDIA, 2026. Kumo tabular: A pretrained tabular foundation model.
\url{https://huggingface.co/nvidia/Kumo-Tabular}

\bibitem[\citeproctext]{ref-operti2017assessing}
Operti, A., Braun, V., Lemmens, S., Corpino, S., 2017. Assessing
uncertainties in the estimation of the orbital lifetime, in: Proceedings
of the 7th European Conference on Space Debris. ESA Space Debris Office,
Darmstadt, Germany.
\url{https://conference.sdo.esoc.esa.int/proceedings/sdc7/paper/1042}

\bibitem[\citeproctext]{ref-pardini2008impact}
Pardini, C., Anselmo, L., 2008. Impact of the time span selected to
calibrate the ballistic parameter on spacecraft re-entry predictions.
Advances in Space Research 41, 1100--1114.
\url{https://doi.org/10.1016/j.asr.2007.11.013}

\bibitem[\citeproctext]{ref-pardini2018assessing}
Pardini, C., Anselmo, L., 2018. Assessing the risk and the uncertainty
affecting the uncontrolled re-entry of manmade space objects. Journal of
Space Safety Engineering 5, 46--62.
\url{https://doi.org/10.1016/j.jsse.2018.01.003}

\bibitem[\citeproctext]{ref-prokhorenkova2018catboost}
Prokhorenkova, L., Gusev, G., Vorobev, A., Dorogush, A.V., Gulin, A.,
2018. {CatBoost}: Unbiased boosting with categorical features, in:
Advances in Neural Information Processing Systems 31.
\url{https://proceedings.neurips.cc/paper/2018/hash/14491b756b3a51daac41c24863285549-Abstract.html}

\bibitem[\citeproctext]{ref-qu2026tabiclv2}
Qu, J., Holzmüller, D., Varoquaux, G., Le Morvan, M., 2026. {TabICLv2}:
A better, faster, scalable, and open tabular foundation model. arXiv
preprint arXiv:2602.11139.
\url{https://doi.org/10.48550/arXiv.2602.11139}

\bibitem[\citeproctext]{ref-romano2019cqr}
Romano, Y., Patterson, E., Candès, E., 2019. Conformalized quantile
regression, in: Advances in Neural Information Processing Systems 32
(NeurIPS 2019). Curran Associates, Inc.
\url{https://proceedings.neurips.cc/paper_files/paper/2019/file/5103c3584b063c431bd1268e9b5e76fb-Paper.pdf}

\bibitem[\citeproctext]{ref-ruszczak2025opssat}
Ruszczak, B., Kotowski, K., Evans, D., Nalepa, J., 2025. The {OPS-SAT}
benchmark for detecting anomalies in satellite telemetry. Sci. Data 12,
710. \url{https://doi.org/10.1038/s41597-025-05035-3}

\bibitem[\citeproctext]{ref-salmaso2023machine}
Salmaso, F., Trisolini, M., Colombo, C., 2023. A machine learning and
feature engineering approach for the prediction of the uncontrolled
re-entry of space objects. Aerospace 10, 297.
\url{https://doi.org/10.3390/aerospace10030297}

\bibitem[\citeproctext]{ref-sang2013estimation}
Sang, J., Bennett, J.C., Smith, C.H., 2013. Estimation of ballistic
coefficients of low altitude debris objects from historical two line
elements. Advances in Space Research 52, 117--124.
\url{https://doi.org/10.1016/j.asr.2013.03.010}

\bibitem[\citeproctext]{ref-shambaugh2026validation}
Shambaugh, S., 2026. Validation of satellite lifetime predictions at
{Leonid} {Space}. \url{https://doi.org/10.48550/arXiv.2601.02453}

\bibitem[\citeproctext]{ref-tulczyjew2021predicting}
Tulczyjew, Ł., Myller, M., Kawulok, M., Kostrzewa, D., Nalepa, J., 2021.
Predicting risk of satellite collisions using machine learning. Journal
of Space Safety Engineering 8, 339--344.
\url{https://doi.org/10.1016/j.jsse.2021.09.001}

\bibitem[\citeproctext]{ref-uriot2021cacdataset}
Uriot, T., Izzo, D., Martinez-Heras, J., Letizia, F., Siminski, J.,
Merz, K., 2021. Collision avoidance challenge dataset.
\url{https://doi.org/10.5281/zenodo.4463683}

\bibitem[\citeproctext]{ref-uriot2022spacecraft}
Uriot, T., Izzo, D., Simões, L.F., Abay, R., Einecke, N., Rebhan, S.,
Martinez-Heras, J., Letizia, F., Siminski, J., Merz, K., 2022.
Spacecraft collision avoidance challenge: Design and results of a
machine learning competition. Astrodynamics 6, 121--140.
\url{https://doi.org/10.1007/s42064-021-0101-5}

\bibitem[\citeproctext]{ref-vazquez2026latest}
Vazquez, N., 2026. Is the latest message still the best forecast? {A}
pre-registered re-evaluation of machine learning for satellite collision
risk. \url{https://doi.org/10.5281/zenodo.23457856}

\bibitem[\citeproctext]{ref-vovk2005algorithmic}
Vovk, V., Gammerman, A., Shafer, G., 2005. Algorithmic learning in a
random world. Springer, New York. \url{https://doi.org/10.1007/b106715}

\bibitem[\citeproctext]{ref-wu2023investigation}
Wu, D., Rosengren, A.J., 2023. An investigation on space debris of
unknown origin using proper elements and neural networks. Celestial
Mechanics and Dynamical Astronomy 135, 44.
\url{https://doi.org/10.1007/s10569-023-10157-0}

\bibitem[\citeproctext]{ref-wu2021flawed}
Wu, R., Keogh, E., 2021. Current time series anomaly detection
benchmarks are flawed and are creating the illusion of progress. IEEE
Trans. Knowl. Data Eng. \url{https://doi.org/10.1109/TKDE.2021.3112126}

\bibitem[\citeproctext]{ref-xu2025digital}
Xu, X., Wen, H., Wang, Y., Song, H., Liu, T., Chang, S.-Y., 2025.
Digital-twin-based satellite orbit prediction for {Internet of Things}
systems. IEEE Internet of Things Journal 12, 6431--6444.
\url{https://doi.org/10.1109/JIOT.2024.3424672}

\bibitem[\citeproctext]{ref-zadrozny2002transforming}
Zadrozny, B., Elkan, C., 2002. Transforming classifier scores into
accurate multiclass probability estimates, in: Proceedings of the Eighth
ACM SIGKDD International Conference on Knowledge Discovery and Data
Mining. ACM, pp. 694--699. \url{https://doi.org/10.1145/775047.775151}

\end{CSLReferences}

\end{document}
